\documentclass[12pt,a4paper]{article}
\usepackage[utf8]{inputenc}
\usepackage[T1]{fontenc}
\usepackage[french]{babel}
\usepackage{amsmath,amssymb}
\usepackage{xcolor}
\usepackage{geometry}
\usepackage{tcolorbox}
\usepackage{tikz}
\usepackage{titlesec}
\usepackage{enumitem}
\usepackage{array}
\usepackage{booktabs}
\usetikzlibrary{arrows.meta, positioning, calc}
\geometry{margin=2cm}

\definecolor{titrebleu}{RGB}{0,102,204}
\definecolor{defvert}{RGB}{0,153,76}
\definecolor{exempleorange}{RGB}{255,102,0}
\definecolor{algoviolet}{RGB}{153,0,153}
\definecolor{importantrouge}{RGB}{204,0,0}
\definecolor{fondgris}{RGB}{245,245,245}
\definecolor{fondvert}{RGB}{230,255,230}
\definecolor{fondbleu}{RGB}{230,240,255}
\definecolor{fondjaune}{RGB}{255,250,205}
\definecolor{fondrose}{RGB}{255,230,240}

\newtcolorbox{enonce}{colback=fondbleu, colframe=titrebleu, fonttitle=\bfseries, title=\'Enonc\'e, arc=3mm}
\newtcolorbox{methode}{colback=fondjaune, colframe=exempleorange, fonttitle=\bfseries, title=M\'ethode, arc=3mm}
\newtcolorbox{correction}{colback=fondvert, colframe=defvert, fonttitle=\bfseries, arc=3mm}
\newtcolorbox{attention}{colback=fondrose, colframe=importantrouge, fonttitle=\bfseries, title=\`A retenir, arc=3mm}
\newtcolorbox{etape}{colback=fondgris, colframe=algoviolet, fonttitle=\bfseries, arc=3mm}

\titleformat{\section}{\color{titrebleu}\normalfont\Large\bfseries}{\color{titrebleu}\thesection}{1em}{}
\titleformat{\subsection}{\color{defvert}\normalfont\large\bfseries}{\color{defvert}\thesubsection}{1em}{}

\title{\textbf{\Huge\color{titrebleu}Exercices Types Corrig\'es}\\[0.5em]
\Large\color{defvert}Optimisation dans les R\'eseaux --- Pr\'eparation Examen\\[0.3em]
\normalsize\color{algoviolet}F.~Bendali-Mailfert --- L3 Informatique --- S6 2025}
\author{}
\date{}

\begin{document}
\maketitle
\tableofcontents
\newpage

%=============================================================
%  EXERCICE 1 : DIJKSTRA
%=============================================================
\section{Exercice 1 --- Dijkstra (poids $\geq 0$)}

\begin{enonce}
Soit le graphe orient\'e $G$ suivant avec source $s$ :

\begin{center}
\begin{tikzpicture}[
    node distance=2.5cm,
    every node/.style={circle, draw, minimum size=9mm, font=\small\bfseries},
    arr/.style={-{Stealth}, thick}
]
    \node[fill=titrebleu!20] (s) at (0,0) {$s$};
    \node[fill=white] (a) at (2.5,1.5) {$a$};
    \node[fill=white] (b) at (2.5,-1.5) {$b$};
    \node[fill=white] (c) at (5,1.5) {$c$};
    \node[fill=white] (d) at (5,-1.5) {$d$};
    \node[fill=importantrouge!20] (t) at (7.5,0) {$t$};

    \draw[arr] (s) -- node[above left, draw=none, fill=none, font=\small]{2} (a);
    \draw[arr] (s) -- node[below left, draw=none, fill=none, font=\small]{5} (b);
    \draw[arr] (a) -- node[above, draw=none, fill=none, font=\small]{1} (c);
    \draw[arr] (a) -- node[right, draw=none, fill=none, font=\small]{4} (d);
    \draw[arr] (b) -- node[below, draw=none, fill=none, font=\small]{2} (d);
    \draw[arr] (b) -- node[left, draw=none, fill=none, font=\small]{1} (a);
    \draw[arr] (c) -- node[above right, draw=none, fill=none, font=\small]{6} (t);
    \draw[arr] (c) -- node[right, draw=none, fill=none, font=\small]{2} (d);
    \draw[arr] (d) -- node[below right, draw=none, fill=none, font=\small]{1} (t);
\end{tikzpicture}
\end{center}

\begin{enumerate}
    \item Quel algorithme choisir ? Justifier.
    \item Appliquer l'algorithme pour trouver les plus courtes distances depuis $s$.
    \item Donner l'arborescence des plus courts chemins.
\end{enumerate}
\end{enonce}

\begin{methode}
\textbf{R\`egle :} Tous les poids $\geq 0$ $\Rightarrow$ \textbf{Dijkstra}.

\textbf{Principe :} \`A chaque \'etape, extraire le sommet non marqu\'e de distance minimale, puis relaxer ses voisins.
\end{methode}

\begin{correction}{title=Correction}

\textbf{Q1.} Tous les poids sont $\geq 0$ $\Rightarrow$ on utilise \textbf{Dijkstra}. Complexit\'e : $O(|E| + |V|\log|V|)$.

\medskip
\textbf{Q2.} Ex\'ecution de Dijkstra :

\begin{center}
\renewcommand{\arraystretch}{1.3}
\begin{tabular}{|c|c|c|c|c|c|c|l|}
\hline
\textbf{\'Etape} & $s$ & $a$ & $b$ & $c$ & $d$ & $t$ & \textbf{Extrait} \\\hline
Init & \textbf{0} & $\infty$ & $\infty$ & $\infty$ & $\infty$ & $\infty$ & --- \\\hline
1 & 0 & \textbf{2} & 5 & $\infty$ & $\infty$ & $\infty$ & $s$ $(d=0)$ \\\hline
2 & 0 & 2 & 5 & \textbf{3} & 6 & $\infty$ & $a$ $(d=2)$ \\\hline
3 & 0 & 2 & 5 & 3 & \textbf{5} & 9 & $c$ $(d=3)$ \\\hline
4 & 0 & 2 & 5 & 3 & 5 & \textbf{6} & $b\text{ ou }d$ $(d=5)$ \\\hline
5 & 0 & 2 & 5 & 3 & 5 & 6 & $d$ $(d=5)$ \\\hline
6 & 0 & 2 & 5 & 3 & 5 & \textbf{6} & $t$ $(d=6)$ \\\hline
\end{tabular}
\end{center}

D\'etail des relaxations :

\begin{etape}{title=\'Etape 1 : extraire $s$ $(d=0)$}
\begin{itemize}[nosep]
    \item $d[a] = \min(\infty,\, 0+2) = 2$ \quad $\pi[a]=s$
    \item $d[b] = \min(\infty,\, 0+5) = 5$ \quad $\pi[b]=s$
\end{itemize}
\end{etape}

\begin{etape}{title=\'Etape 2 : extraire $a$ $(d=2)$}
\begin{itemize}[nosep]
    \item $d[c] = \min(\infty,\, 2+1) = 3$ \quad $\pi[c]=a$
    \item $d[d] = \min(\infty,\, 2+4) = 6$ \quad $\pi[d]=a$
\end{itemize}
\end{etape}

\begin{etape}{title=\'Etape 3 : extraire $c$ $(d=3)$}
\begin{itemize}[nosep]
    \item $d[t] = \min(\infty,\, 3+6) = 9$ \quad $\pi[t]=c$
    \item $d[d] = \min(6,\, 3+2) = \mathbf{5}$ \quad $\pi[d]=c$ \quad (am\'elior\'e !)
\end{itemize}
\end{etape}

\begin{etape}{title=\'Etape 4 : extraire $b$ $(d=5)$}
\begin{itemize}[nosep]
    \item $d[d] = \min(5,\, 5+2) = 5$ \quad (inchang\'e)
    \item $d[a] = \min(2,\, 5+1) = 2$ \quad (inchang\'e, d\'ej\`a marqu\'e)
\end{itemize}
\end{etape}

\begin{etape}{title=\'Etape 5 : extraire $d$ $(d=5)$}
\begin{itemize}[nosep]
    \item $d[t] = \min(9,\, 5+1) = \mathbf{6}$ \quad $\pi[t]=d$ \quad (am\'elior\'e !)
\end{itemize}
\end{etape}

\medskip
\textbf{Q3.} Arborescence des plus courts chemins :

$$\pi[a]=s,\quad \pi[b]=s,\quad \pi[c]=a,\quad \pi[d]=c,\quad \pi[t]=d$$

\begin{center}
\begin{tikzpicture}[
    every node/.style={circle, draw, minimum size=9mm, font=\small\bfseries},
    arr/.style={-{Stealth}, thick, color=defvert}
]
    \node[fill=titrebleu!20] (s) at (0,0) {$s$};
    \node (a) at (2,1) {$a$};
    \node (b) at (2,-1) {$b$};
    \node (c) at (4,1) {$c$};
    \node (d) at (6,0) {$d$};
    \node[fill=importantrouge!20] (t) at (8,0) {$t$};

    \draw[arr] (s) -- node[above left, draw=none, fill=none, font=\scriptsize, color=black]{$d\!=\!2$} (a);
    \draw[arr] (s) -- node[below left, draw=none, fill=none, font=\scriptsize, color=black]{$d\!=\!5$} (b);
    \draw[arr] (a) -- node[above, draw=none, fill=none, font=\scriptsize, color=black]{$d\!=\!3$} (c);
    \draw[arr] (c) -- node[above, draw=none, fill=none, font=\scriptsize, color=black]{$d\!=\!5$} (d);
    \draw[arr] (d) -- node[above, draw=none, fill=none, font=\scriptsize, color=black]{$d\!=\!6$} (t);
\end{tikzpicture}
\end{center}

\textbf{PCC de $s$ \`a $t$ :} $s \to a \to c \to d \to t$ \quad (co\^ut = $2+1+2+1 = 6$)
\end{correction}

\begin{attention}
\begin{itemize}
    \item On extrait toujours le sommet \textbf{non marqu\'e de distance minimale}.
    \item Un sommet d\'ej\`a marqu\'e ne change plus de distance.
    \item Dijkstra \textbf{ne marche pas} si poids n\'egatifs $\Rightarrow$ utiliser Bellman-Ford ou relaxation topologique (DAG).
\end{itemize}
\end{attention}


\newpage
%=============================================================
%  EXERCICE 2 : RELAXATION TOPOLOGIQUE (DAG + POIDS NEGATIFS)
%=============================================================
\section{Exercice 2 --- Relaxation topologique (DAG, poids n\'egatifs)}

\begin{enonce}
Soit le graphe orient\'e $G$ suivant. L'arc $a \to c$ a un poids \textbf{n\'egatif} $(-4)$.

\begin{center}
\begin{tikzpicture}[
    every node/.style={circle, draw, minimum size=9mm, font=\small\bfseries},
    arr/.style={-{Stealth}, thick}
]
    \node[fill=titrebleu!20] (s) at (0,0) {$s$};
    \node (a) at (2.5,1.5) {$a$};
    \node (b) at (2.5,-1.5) {$b$};
    \node (c) at (5,0) {$c$};
    \node (d) at (7.5,1.5) {$d$};
    \node[fill=importantrouge!20] (e) at (7.5,-1.5) {$e$};

    \draw[arr] (s) -- node[above left, draw=none, fill=none, font=\small]{3} (a);
    \draw[arr] (s) -- node[below left, draw=none, fill=none, font=\small]{2} (b);
    \draw[arr] (a) -- node[above, draw=none, fill=none, font=\small, color=importantrouge]{\textbf{-4}} (c);
    \draw[arr] (a) -- node[above right, draw=none, fill=none, font=\small]{5} (d);
    \draw[arr] (b) -- node[below, draw=none, fill=none, font=\small]{3} (c);
    \draw[arr] (b) -- node[below right, draw=none, fill=none, font=\small]{7} (e);
    \draw[arr] (c) -- node[right, draw=none, fill=none, font=\small]{2} (d);
    \draw[arr] (c) -- node[right, draw=none, fill=none, font=\small]{6} (e);
    \draw[arr] (d) -- node[right, draw=none, fill=none, font=\small]{1} (e);
\end{tikzpicture}
\end{center}

\begin{enumerate}
    \item Montrer que $G$ est un DAG en donnant les niveaux.
    \item Quel algorithme choisir ? Pourquoi pas Dijkstra ?
    \item Calculer les PCC depuis $s$ par relaxation topologique.
    \item Donner le PCC de $s$ \`a $e$.
\end{enumerate}
\end{enonce}

\begin{methode}
\textbf{R\`egle :} Poids n\'egatifs $\Rightarrow$ \textbf{pas Dijkstra}. Si le graphe est un DAG $\Rightarrow$ \textbf{relaxation topologique} (meilleur que Bellman-Ford : $O(V+E)$).

\textbf{Pour trouver les niveaux :} Chercher les sommets sans pr\'ed\'ecesseur (niveau 0), puis it\'erer.
\end{methode}

\begin{correction}{title=Correction}

\textbf{Q1. Niveaux :}

Pr\'ed\'ecesseurs de chaque sommet :
\begin{itemize}[nosep]
    \item $s$ : aucun $\Rightarrow$ \textbf{Niveau 0}
    \item $a$ : $s$ (niv.0) $\Rightarrow$ \textbf{Niveau 1}
    \item $b$ : $s$ (niv.0) $\Rightarrow$ \textbf{Niveau 1}
    \item $c$ : $a$ (niv.1), $b$ (niv.1) $\Rightarrow$ \textbf{Niveau 2}
    \item $d$ : $a$ (niv.1), $c$ (niv.2) $\Rightarrow$ \textbf{Niveau 3}
    \item $e$ : $b$ (niv.1), $c$ (niv.2), $d$ (niv.3) $\Rightarrow$ \textbf{Niveau 4}
\end{itemize}

Tous les arcs vont d'un niveau inf\'erieur vers un niveau sup\'erieur $\Rightarrow$ \textbf{pas de circuit} $\Rightarrow$ DAG. $\checkmark$

\medskip
\textbf{Q2.} L'arc $a \to c$ a le poids $-4 < 0$ $\Rightarrow$ Dijkstra \textbf{non applicable}. Mais $G$ est un DAG $\Rightarrow$ on utilise la \textbf{relaxation par ordre topologique} en $O(V+E)$.

\medskip
\textbf{Q3. Ex\'ecution :}

Init : $d[s]=0$, tous les autres $= +\infty$.

\begin{etape}{title=Niveau 0 --- traiter $s$}
$d[a] = \min(\infty,\, 0+3) = 3$, \quad $d[b] = \min(\infty,\, 0+2) = 2$
\end{etape}

\begin{etape}{title=Niveau 1 --- traiter $a$ ($d=3$) puis $b$ ($d=2$)}
Depuis $a$ : $d[c] = \min(\infty,\, 3+(-4)) = \mathbf{-1}$, \quad $d[d] = \min(\infty,\, 3+5) = 8$

Depuis $b$ : $d[c] = \min(-1,\, 2+3) = -1$ (inchang\'e), \quad $d[e] = \min(\infty,\, 2+7) = 9$
\end{etape}

\begin{etape}{title=Niveau 2 --- traiter $c$ ($d=-1$)}
$d[d] = \min(8,\, -1+2) = \mathbf{1}$ \quad (am\'elior\'e !), \quad $d[e] = \min(9,\, -1+6) = \mathbf{5}$ \quad (am\'elior\'e !)
\end{etape}

\begin{etape}{title=Niveau 3 --- traiter $d$ ($d=1$)}
$d[e] = \min(5,\, 1+1) = \mathbf{2}$ \quad (am\'elior\'e !)
\end{etape}

\begin{center}
\renewcommand{\arraystretch}{1.3}
\begin{tabular}{|c|c|c|c|c|c|c|}
\hline
\textbf{Sommet} & $s$ & $a$ & $b$ & $c$ & $d$ & $e$ \\\hline
$d(\cdot)$ & 0 & 3 & 2 & $-1$ & 1 & 2 \\\hline
$\pi(\cdot)$ & --- & $s$ & $s$ & $a$ & $c$ & $d$ \\\hline
\end{tabular}
\end{center}

\textbf{Q4. PCC de $s$ \`a $e$ :} $s \to a \to c \to d \to e$ \quad (co\^ut = $3 + (-4) + 2 + 1 = \mathbf{2}$)
\end{correction}

\begin{attention}
\begin{itemize}
    \item Le poids n\'egatif sur $a \to c$ permet d'obtenir $d[c] = -1$ : c'est normal dans un DAG (pas de circuit absorbant possible).
    \item L'ordre de traitement par niveaux garantit qu'on ne revient jamais en arri\`ere.
\end{itemize}
\end{attention}


\newpage
%=============================================================
%  EXERCICE 3 : KRUSKAL + PUISSANCE
%=============================================================
\section{Exercice 3 --- Kruskal + Puissance des n\oe uds}

\begin{enonce}
Soit le graphe non orient\'e $G$ (r\'eseau sans fil) :

\begin{center}
\begin{tikzpicture}[
    every node/.style={circle, draw, minimum size=9mm, font=\small\bfseries},
    edge/.style={thick}
]
    \node (1) at (0,2) {1};
    \node (2) at (2,3) {2};
    \node (3) at (4,2) {3};
    \node (4) at (2,0) {4};
    \node (5) at (6,1) {5};
    \node (6) at (4,0) {6};

    \draw[edge] (1) -- node[above, draw=none, fill=none, font=\small]{3} (2);
    \draw[edge] (1) -- node[left, draw=none, fill=none, font=\small]{5} (4);
    \draw[edge] (2) -- node[above, draw=none, fill=none, font=\small]{4} (3);
    \draw[edge] (2) -- node[right, draw=none, fill=none, font=\small]{6} (4);
    \draw[edge] (3) -- node[right, draw=none, fill=none, font=\small]{2} (5);
    \draw[edge] (3) -- node[right, draw=none, fill=none, font=\small]{7} (6);
    \draw[edge] (4) -- node[below, draw=none, fill=none, font=\small]{8} (6);
    \draw[edge] (5) -- node[right, draw=none, fill=none, font=\small]{5} (6);
    \draw[edge] (1) -- node[above right, draw=none, fill=none, font=\small]{9} (3);
\end{tikzpicture}
\end{center}

\begin{enumerate}
    \item Appliquer Kruskal pour trouver un MST.
    \item Calculer la puissance de chaque n\oe ud dans le MST.
    \item Calculer la puissance totale du r\'eseau.
\end{enumerate}
\end{enonce}

\begin{methode}
\textbf{Kruskal :} Trier les ar\^etes par poids croissant. Ajouter si pas de cycle. S'arr\^eter \`a $|V|-1$ ar\^etes.

\textbf{Puissance :} $p(v) = \max\{$poids des ar\^etes incidentes \`a $v$ dans le MST$\}$.
\end{methode}

\begin{correction}{title=Correction}

\textbf{Q1. Kruskal :}

Ar\^etes tri\'ees : $3$-$5$(2), $1$-$2$(3), $2$-$3$(4), $1$-$4$(5), $5$-$6$(5), $2$-$4$(6), $3$-$6$(7), $4$-$6$(8), $1$-$3$(9)

\begin{center}
\renewcommand{\arraystretch}{1.3}
\begin{tabular}{|c|c|c|c|}
\hline
\textbf{\'Etape} & \textbf{Ar\^ete} & \textbf{Poids} & \textbf{Action} \\\hline
1 & $3$-$5$ & 2 & Ajouter $\to \{3,5\}$ \\\hline
2 & $1$-$2$ & 3 & Ajouter $\to \{1,2\},\{3,5\}$ \\\hline
3 & $2$-$3$ & 4 & Ajouter $\to \{1,2,3,5\}$ \\\hline
4 & $1$-$4$ & 5 & Ajouter $\to \{1,2,3,4,5\}$ \\\hline
5 & $5$-$6$ & 5 & Ajouter $\to \{1,2,3,4,5,6\}$ \checkmark \\\hline
\end{tabular}
\end{center}

$|V|-1 = 5$ ar\^etes s\'electionn\'ees. \textbf{Co\^ut MST} $= 2+3+4+5+5 = \mathbf{19}$.

\begin{center}
\begin{tikzpicture}[
    every node/.style={circle, draw, minimum size=9mm, font=\small\bfseries},
    edge/.style={thick, color=defvert},
    unused/.style={thick, dashed, color=gray!40}
]
    \node (1) at (0,2) {1};
    \node (2) at (2,3) {2};
    \node (3) at (4,2) {3};
    \node (4) at (2,0) {4};
    \node (5) at (6,1) {5};
    \node (6) at (4,0) {6};

    \draw[edge] (1) -- node[above, draw=none, fill=none, font=\small]{3} (2);
    \draw[edge] (2) -- node[above, draw=none, fill=none, font=\small]{4} (3);
    \draw[edge] (1) -- node[left, draw=none, fill=none, font=\small]{5} (4);
    \draw[edge] (3) -- node[right, draw=none, fill=none, font=\small]{2} (5);
    \draw[edge] (5) -- node[right, draw=none, fill=none, font=\small]{5} (6);

    \draw[unused] (2) -- (4);
    \draw[unused] (3) -- (6);
    \draw[unused] (4) -- (6);
    \draw[unused] (1) -- (3);
\end{tikzpicture}
\end{center}

\medskip
\textbf{Q2--Q3. Puissance de chaque n\oe ud :}

\begin{center}
\renewcommand{\arraystretch}{1.3}
\begin{tabular}{|c|l|c|}
\hline
\textbf{N\oe ud} & \textbf{Ar\^etes incidentes dans MST} & $p(v)$ \\\hline
1 & $1$-$2$(3), $1$-$4$(5) & $\max(3,5) = \mathbf{5}$ \\\hline
2 & $1$-$2$(3), $2$-$3$(4) & $\max(3,4) = \mathbf{4}$ \\\hline
3 & $2$-$3$(4), $3$-$5$(2) & $\max(4,2) = \mathbf{4}$ \\\hline
4 & $1$-$4$(5) & $\mathbf{5}$ \\\hline
5 & $3$-$5$(2), $5$-$6$(5) & $\max(2,5) = \mathbf{5}$ \\\hline
6 & $5$-$6$(5) & $\mathbf{5}$ \\\hline\hline
\multicolumn{2}{|r|}{\textbf{Puissance totale}} & $\mathbf{28}$ \\\hline
\end{tabular}
\end{center}
\end{correction}

\begin{attention}
\begin{itemize}
    \item $p(v) = \max$ (pas la somme !) des poids des ar\^etes incidentes \`a $v$ dans l'arbre.
    \item Si deux ar\^etes ont le m\^eme poids, il peut y avoir \textbf{plusieurs MST} de m\^eme co\^ut mais de puissances diff\'erentes.
\end{itemize}
\end{attention}


\newpage
%=============================================================
%  EXERCICE 4 : BIP
%=============================================================
\section{Exercice 4 --- Algorithme BIP}

\begin{enonce}
Sur le m\^eme graphe $G$ de l'exercice~3, appliquer l'algorithme \textbf{BIP} (Broadcast Incremental Power) depuis le n\oe ud~$1$.

Comparer la puissance totale obtenue avec celle du MST.
\end{enonce}

\begin{methode}
\textbf{BIP :} \`A chaque \'etape, choisir l'ar\^ete $(u,v)$ ($u$ dans l'arbre, $v$ hors) qui minimise le \textbf{co\^ut additionnel} :
$$\text{co\^ut\_add}(u,v) = \max\big(C(u,v) - p(u),\; 0\big)$$
Si $C(u,v) \leq p(u)$ : le co\^ut additionnel est $0$ (le n\oe ud \'emet d\'ej\`a assez fort).
\end{methode}

\begin{correction}{title=Correction}

\begin{etape}{title=\'Etape 1 : Arbre $= \{1\}$, $p(1)=0$}
Ar\^etes depuis 1 :
\begin{itemize}[nosep]
    \item $1$-$2$(3) : co\^ut\_add $= 3-0 = 3$
    \item $1$-$4$(5) : co\^ut\_add $= 5-0 = 5$
    \item $1$-$3$(9) : co\^ut\_add $= 9-0 = 9$
\end{itemize}
\textbf{Min :} $1$-$2$(3), co\^ut add $= 3$.  $\Rightarrow p(1) = 3$, $p(2) = 3$.
\end{etape}

\begin{etape}{title=\'Etape 2 : Arbre $= \{1,2\}$, $p(1)=3$, $p(2)=3$}
\begin{itemize}[nosep]
    \item $1$-$4$(5) : $\max(5-3,0) = 2$
    \item $1$-$3$(9) : $\max(9-3,0) = 6$
    \item $2$-$3$(4) : $\max(4-3,0) = 1$
    \item $2$-$4$(6) : $\max(6-3,0) = 3$
\end{itemize}
\textbf{Min :} $2$-$3$(4), co\^ut add $= 1$. $\Rightarrow p(2)=4$, $p(3)=4$.
\end{etape}

\begin{etape}{title=\'Etape 3 : Arbre $= \{1,2,3\}$, $p(1)=3$, $p(2)=4$, $p(3)=4$}
\begin{itemize}[nosep]
    \item $1$-$4$(5) : $\max(5-3,0) = 2$
    \item $3$-$5$(2) : $\max(2-4,0) = \mathbf{0}$ \quad (gratuit !)
    \item $3$-$6$(7) : $\max(7-4,0) = 3$
\end{itemize}
\textbf{Min :} $3$-$5$(2), co\^ut add $= 0$. $\Rightarrow p(3)$ inchang\'e (4), $p(5)=2$.
\end{etape}

\begin{etape}{title=\'Etape 4 : Arbre $= \{1,2,3,5\}$}
\begin{itemize}[nosep]
    \item $1$-$4$(5) : $\max(5-3,0) = 2$
    \item $5$-$6$(5) : $\max(5-2,0) = 3$
    \item $3$-$6$(7) : $\max(7-4,0) = 3$
\end{itemize}
\textbf{Min :} $1$-$4$(5), co\^ut add $= 2$. $\Rightarrow p(1)=5$, $p(4)=5$.
\end{etape}

\begin{etape}{title=\'Etape 5 : Arbre $= \{1,2,3,4,5\}$}
\begin{itemize}[nosep]
    \item $5$-$6$(5) : $\max(5-2,0) = 3$
    \item $3$-$6$(7) : $\max(7-4,0) = 3$
    \item $4$-$6$(8) : $\max(8-5,0) = 3$
\end{itemize}
Ex-\ae quo. Choisir $5$-$6$(5). $\Rightarrow p(5)=5$, $p(6)=5$.
\end{etape}

\medskip
\textbf{R\'esultat BIP :}

\begin{center}
\renewcommand{\arraystretch}{1.3}
\begin{tabular}{|c|c|c|c|c|c|c||c|}
\hline
\textbf{N\oe ud} & 1 & 2 & 3 & 4 & 5 & 6 & \textbf{Total} \\\hline
$p_{\text{BIP}}$ & 5 & 4 & 4 & 5 & 5 & 5 & $\mathbf{28}$ \\\hline
$p_{\text{MST}}$ & 5 & 4 & 4 & 5 & 5 & 5 & $\mathbf{28}$ \\\hline
\end{tabular}
\end{center}

Ici les deux m\'ethodes donnent la m\^eme puissance. Ce n'est \textbf{pas toujours le cas} --- voir l'examen o\`u BIP donne 44 vs MST $T_1$ = 47.
\end{correction}

\begin{attention}
\begin{itemize}
    \item Co\^ut additionnel $= \max(C(u,v) - p(u),\, 0)$ : ne jamais oublier le $\max(\cdot, 0)$.
    \item Un co\^ut additionnel de $0$ signifie que $u$ \'emet d\'ej\`a assez fort pour atteindre $v$ ``gratuitement''.
    \item BIP \textbf{ne garantit pas l'optimum} global, c'est une heuristique gloutonne.
\end{itemize}
\end{attention}


\newpage
%=============================================================
%  EXERCICE 5 : FLOT MAX + COUPE MIN + CHEMINS DISJOINTS
%=============================================================
\section{Exercice 5 --- Flot maximum, coupe minimum, chemins arcs-disjoints}

\begin{enonce}
Soit le r\'eseau $G$ suivant avec capacit\'es sur les arcs :

\begin{center}
\begin{tikzpicture}[
    every node/.style={circle, draw, minimum size=9mm, font=\small\bfseries},
    arr/.style={-{Stealth}, thick}
]
    \node[fill=titrebleu!20] (s) at (0,0) {$s$};
    \node (a) at (2.5,2) {$a$};
    \node (b) at (2.5,0) {$b$};
    \node (c) at (2.5,-2) {$c$};
    \node (d) at (5.5,1) {$d$};
    \node (e) at (5.5,-1) {$e$};
    \node[fill=importantrouge!20] (t) at (8,0) {$t$};

    \draw[arr] (s) -- node[above left, draw=none, fill=none, font=\small]{5} (a);
    \draw[arr] (s) -- node[above, draw=none, fill=none, font=\small]{4} (b);
    \draw[arr] (s) -- node[below left, draw=none, fill=none, font=\small]{3} (c);
    \draw[arr] (a) -- node[above, draw=none, fill=none, font=\small]{3} (d);
    \draw[arr] (b) -- node[above, draw=none, fill=none, font=\small]{2} (d);
    \draw[arr] (b) -- node[below, draw=none, fill=none, font=\small]{3} (e);
    \draw[arr] (c) -- node[below, draw=none, fill=none, font=\small]{4} (e);
    \draw[arr] (d) -- node[above right, draw=none, fill=none, font=\small]{6} (t);
    \draw[arr] (e) -- node[below right, draw=none, fill=none, font=\small]{5} (t);
\end{tikzpicture}
\end{center}

\begin{enumerate}
    \item Trouver le flot maximum de $s$ \`a $t$ en utilisant Ford-Fulkerson (chemins augmentants).
    \item Donner une coupe minimum et v\'erifier Flot-Max $=$ Coupe-Min.
    \item En mettant les capacit\'es \`a 1, combien de chemins arcs-disjoints $s$-$t$ au maximum ?
\end{enumerate}
\end{enonce}

\begin{methode}
\textbf{Ford-Fulkerson :}
\begin{enumerate}[nosep]
    \item Partir de flot $= 0$
    \item Trouver un chemin augmentant de $s$ \`a $t$ dans le graphe r\'esiduel
    \item Augmenter le flot du min des capacit\'es r\'esiduelles sur le chemin
    \item R\'ep\'eter jusqu'\`a plus de chemin augmentant
\end{enumerate}
\end{methode}

\begin{correction}{title=Correction}

\textbf{Q1. Ford-Fulkerson :}

\begin{etape}{title=Chemin augmentant 1 : $s \to a \to d \to t$}
Capacit\'es r\'esiduelles : $\min(5, 3, 6) = 3$. Envoyer $3$ unit\'es.

Flots : $x_{sa}=3$, $x_{ad}=3$, $x_{dt}=3$. Valeur du flot : $\mathbf{3}$.
\end{etape}

\begin{etape}{title=Chemin augmentant 2 : $s \to b \to e \to t$}
Capacit\'es r\'esiduelles : $\min(4, 3, 5) = 3$. Envoyer $3$ unit\'es.

Flots : $x_{sb}=3$, $x_{be}=3$, $x_{et}=3$. Valeur du flot : $\mathbf{6}$.
\end{etape}

\begin{etape}{title=Chemin augmentant 3 : $s \to c \to e \to t$}
Capacit\'es r\'esiduelles : $\min(3, 4, 5-3) = \min(3, 4, 2) = 2$. Envoyer $2$ unit\'es.

Flots : $x_{sc}=2$, $x_{ce}=2$, $x_{et}=5$. Valeur du flot : $\mathbf{8}$.
\end{etape}

\begin{etape}{title=Chemin augmentant 4 : $s \to b \to d \to t$}
Capacit\'es r\'esiduelles : $\min(4-3, 2, 6-3) = \min(1, 2, 3) = 1$. Envoyer $1$ unit\'e.

Flots : $x_{sb}=4$, $x_{bd}=1$, $x_{dt}=4$. Valeur du flot : $\mathbf{9}$.
\end{etape}

\begin{etape}{title=Chemin augmentant 5 : $s \to a \to d \to t$}
R\'esiduel : $\min(5-3, 3-3, \ldots) = \min(2, 0, \ldots) = 0$ sur $a \to d$. \textbf{Bloqu\'e.}

$s \to c \to e \to t$ : $\min(3-2, 4-2, 5-5) = \min(1, 2, 0) = 0$. \textbf{Bloqu\'e.}

Plus de chemin augmentant. $\Rightarrow$ \textbf{Flot maximum $= 9$}.
\end{etape}

\medskip
Flots finaux :

\begin{center}
\begin{tabular}{|c|c|c|}
\hline
\textbf{Arc} & \textbf{Flot / Capacit\'e} & \textbf{Satur\'e ?} \\\hline
$s \to a$ & 3/5 & Non \\\hline
$s \to b$ & 4/4 & \textbf{OUI} \\\hline
$s \to c$ & 2/3 & Non \\\hline
$a \to d$ & 3/3 & \textbf{OUI} \\\hline
$b \to d$ & 1/2 & Non \\\hline
$b \to e$ & 3/3 & \textbf{OUI} \\\hline
$c \to e$ & 2/4 & Non \\\hline
$d \to t$ & 4/6 & Non \\\hline
$e \to t$ & 5/5 & \textbf{OUI} \\\hline
\end{tabular}
\end{center}

V\'erification : sortie de $s$ = $3+4+2 = 9$, entr\'ee dans $t$ = $4+5 = 9$. $\checkmark$

\medskip
\textbf{Q2. Coupe minimum :}

Sommets atteignables depuis $s$ dans le graphe r\'esiduel : $S = \{s, a, c\}$.

$T = \{b, d, e, t\}$.

Arcs de $S$ vers $T$ :
\begin{itemize}[nosep]
    \item $s \to b$ : capacit\'e $4$
    \item $a \to d$ : capacit\'e $3$
    \item $c \to e$ : capacit\'e $4$ \quad ... attendons.
\end{itemize}

En fait, cherchons mieux. L'ensemble $S = \{s\}$, $T = V \setminus \{s\}$ :
$$\text{cap}(\{s\}, T) = 5 + 4 + 3 = 12 \neq 9$$

Essayons $S = \{s, a, c\}$ : arcs $s \to b$ (4), $a \to d$ (3), $c \to e$ (4) $\Rightarrow$ cap $= 11 \neq 9$.

Essayons $S = \{s, a, b, c\}$ : arcs $a \to d$ (3), $b \to d$ (2), $b \to e$ (3), $c \to e$ (4) $\Rightarrow$ cap $= 12 \neq 9$.

Essayons $S = \{s, a, b, c, d\}$ : arcs $b \to e$ (3), $c \to e$ (4), $d \to t$ (6) $\Rightarrow$ cap $= 13$.

Essayons $S = \{s, a, b, c, e\}$ : arcs $a \to d$ (3), $b \to d$ (2), $e \to t$ (5) $\Rightarrow$ cap $= 10$.

Essayons $S = \{s, a, b, c, d, e\}$ : arcs $d \to t$ (6), $e \to t$ (5) $\Rightarrow$ cap $= 11$.

\textbf{Coupe correcte :} $S = \{s, c\}$, arcs $s \to a$ (5), $s \to b$ (4), $c \to e$ (4) $\Rightarrow$ cap $= 13$.

Reprenons le graphe r\'esiduel apr\`es flot max. Les sommets atteignables depuis $s$ :
\begin{itemize}[nosep]
    \item $s$ : directement. $s \to a$ r\'esiduel $5-3=2 > 0 \Rightarrow a$ atteignable.
    \item $a$ : $a \to d$ r\'esiduel $3-3=0$, bloqu\'e. Pas de sortie.
    \item $s \to b$ : r\'esiduel $4-4=0$. Bloqu\'e.
    \item $s \to c$ : r\'esiduel $3-2=1 > 0 \Rightarrow c$ atteignable.
    \item $c \to e$ : r\'esiduel $4-2=2 > 0 \Rightarrow e$ atteignable.
    \item $e \to t$ : r\'esiduel $5-5=0$. Bloqu\'e. Arc retour $e \to b$ : $x_{be}=3 > 0 \Rightarrow b$ atteignable.
    \item $b \to d$ : r\'esiduel $2-1=1 > 0 \Rightarrow d$ atteignable.
    \item $d \to t$ : r\'esiduel $6-4=2 > 0 \Rightarrow t$ atteignable ?!
\end{itemize}

Il reste un chemin ! Chemin : $s \to c \to e \xrightarrow{\text{retour}} b \to d \to t$, avec capacit\'es r\'esiduelles $\min(1, 2, 3, 1, 2) = 1$.

\medskip
\begin{etape}{title=Chemin augmentant 5 : $s \to c \to e \xrightarrow{\text{retour}} b \to d \to t$ (via arc retour)}
$\delta = 1$. Envoyer 1 unit\'e.

Mise \`a jour :
$x_{sc} = 3$, $x_{ce} = 3$, $x_{be} = 2$ (diminu\'e !), $x_{bd} = 2$, $x_{dt} = 5$.

Valeur du flot : $\mathbf{10}$.
\end{etape}

Recontr\^ole : sortie de $s$ = $3+4+3 = 10$, entr\'ee dans $t$ = $5+5 = 10$. $\checkmark$

V\'erifions s'il reste un chemin augmentant :
\begin{itemize}[nosep]
    \item $s \to a$ (r\'es. 2) $\to d$ (r\'es. 0) : bloqu\'e.
    \item $s \to c$ (r\'es. 0) : bloqu\'e.
    \item $s \to b$ (r\'es. 0) : bloqu\'e.
\end{itemize}
Seul $s \to a$ est possible, mais $a$ ne m\`ene nulle part. \textbf{Plus de chemin augmentant.}

\textbf{Flot maximum $= 10$.}

$S = \{s, a\}$ (sommets atteignables). $T = \{b, c, d, e, t\}$.

Arcs de $S$ vers $T$ : $s \to b$ (cap. 4), $s \to c$ (cap. 3), $a \to d$ (cap. 3) $\Rightarrow$ cap $= 4+3+3 = \mathbf{10}$ $= f^*$. $\checkmark$

\textbf{Coupe minimum :} $\{s \to b,\; s \to c,\; a \to d\}$, capacit\'e $= 10 = $ Flot Max. \checkmark

\medskip
\textbf{Q3. Chemins arcs-disjoints (capacit\'es unitaires) :}

Avec capacit\'e $= 1$ sur chaque arc :

On identifie 3 chemins arcs-disjoints :
\begin{itemize}
    \item $P_1 : s \to a \to d \to t$
    \item $P_2 : s \to b \to e \to t$
    \item $P_3 : s \to c \to e \to t$ \quad \textcolor{importantrouge}{\textbf{IMPOSSIBLE}} : l'arc $e \to t$ est d\'ej\`a utilis\'e dans $P_2$.
\end{itemize}

Reconsid\'erons :
\begin{itemize}
    \item $P_1 : s \to a \to d \to t$
    \item $P_2 : s \to b \to e \to t$
    \item $P_3$ n\'ecessiterait $s \to c \to e$, mais $e \to t$ d\'ej\`a pris. Et $b \to d$ : $d \to t$ d\'ej\`a pris.
\end{itemize}

Un 3\textsuperscript{e} chemin ne peut pas \'eviter de partager un arc. Donc flot max (cap. unit.) $= 2$.

V\'erifions : coupe $\{d \to t, e \to t\}$, capacit\'e $= 1+1 = 2$. Confirme que max $= \mathbf{2}$ chemins arcs-disjoints.
\end{correction}

\begin{attention}
\begin{itemize}
    \item \textbf{Ne pas oublier les arcs retour} dans le graphe r\'esiduel ! Un arc $(u,v)$ avec flot $x > 0$ cr\'ee un arc retour $(v,u)$ de capacit\'e $x$.
    \item \textbf{Coupe min :} Sommets atteignables depuis $s$ dans le r\'esiduel $= S$. Arcs $S \to V \setminus S$ $=$ coupe.
    \item \textbf{Chemins arcs-disjoints :} mettre toutes les capacit\'es \`a 1 et r\'esoudre le flot max.
\end{itemize}
\end{attention}


\newpage
%=============================================================
%  EXERCICE 6 : DEDOUBLEMENT DE NOEUD
%=============================================================
\section{Exercice 6 --- D\'edoublement de n\oe ud (capacit\'e nodale)}

\begin{enonce}
Reprendre le graphe de l'exercice 5 avec les capacit\'es originales. On impose maintenant une \textbf{contrainte de capacit\'e sur le n\oe ud $b$} : au plus $2$ unit\'es de flot peuvent transiter par $b$.

\begin{enumerate}
    \item Comment mod\'eliser cette contrainte ?
    \item Quel est le nouveau flot maximum ?
\end{enumerate}
\end{enonce}

\begin{methode}
\textbf{D\'edoublement de n\oe ud :} Remplacer $b$ par $b_{in}$ et $b_{out}$ reli\'es par un arc de capacit\'e \'egale \`a la capacit\'e nodale.
\begin{itemize}[nosep]
    \item Arcs entrant en $b$ $\to$ arrivent \`a $b_{in}$
    \item Arcs sortant de $b$ $\to$ partent de $b_{out}$
    \item Arc $b_{in} \to b_{out}$ de capacit\'e $= 2$
\end{itemize}
\end{methode}

\begin{correction}{title=Correction}

\textbf{Q1. Transformation :}

\begin{center}
\begin{tikzpicture}[
    every node/.style={circle, draw, minimum size=9mm, font=\small\bfseries},
    arr/.style={-{Stealth}, thick}
]
    \node[fill=titrebleu!20] (s) at (0,0) {$s$};
    \node (a) at (2.5,2) {$a$};
    \node[fill=fondjaune] (bin) at (2,-0.5) {$b_{\text{in}}$};
    \node[fill=fondjaune] (bout) at (3.8,-0.5) {$b_{\text{out}}$};
    \node (c) at (2.5,-2.5) {$c$};
    \node (d) at (6,1) {$d$};
    \node (e) at (6,-1.5) {$e$};
    \node[fill=importantrouge!20] (t) at (8.5,0) {$t$};

    \draw[arr] (s) -- node[above left, draw=none, fill=none, font=\small]{5} (a);
    \draw[arr] (s) -- node[above, draw=none, fill=none, font=\small]{4} (bin);
    \draw[arr] (s) -- node[below left, draw=none, fill=none, font=\small]{3} (c);
    \draw[arr, color=importantrouge, very thick] (bin) -- node[above, draw=none, fill=none, font=\small, color=importantrouge]{\textbf{2}} (bout);
    \draw[arr] (a) -- node[above, draw=none, fill=none, font=\small]{3} (d);
    \draw[arr] (bout) -- node[above, draw=none, fill=none, font=\small]{2} (d);
    \draw[arr] (bout) -- node[below, draw=none, fill=none, font=\small]{3} (e);
    \draw[arr] (c) -- node[below, draw=none, fill=none, font=\small]{4} (e);
    \draw[arr] (d) -- node[above right, draw=none, fill=none, font=\small]{6} (t);
    \draw[arr] (e) -- node[below right, draw=none, fill=none, font=\small]{5} (t);
\end{tikzpicture}
\end{center}

\textbf{Q2. Nouveau flot max :}

Avec la contrainte $b_{in} \to b_{out}$ : capacit\'e $2$, le flot passant par $b$ est limit\'e \`a $2$.

Chemins augmentants :
\begin{itemize}[nosep]
    \item $s \to a \to d \to t$ : flot $= 3$ (capacit\'e de $a \to d$)
    \item $s \to b_{in} \to b_{out} \to e \to t$ : flot $= 2$ (goulot : arc $b_{in} \to b_{out}$)
    \item $s \to c \to e \to t$ : r\'esiduel $e \to t$ : $5-2=3$, et $c \to e$ : $4$, $s \to c$ : $3$ $\Rightarrow \min(3,4,3) = 3$.
\end{itemize}

Flot total $= 3 + 2 + 3 = 8$. V\'erifions s'il reste un chemin augmentant.

Apr\`es ces 3 envois : arc $b_{in} \to b_{out}$ satur\'e, $e \to t$ satur\'e ($2+3=5$), $s \to c$ satur\'e.

R\'esiduel depuis $s$ : $s \to a$ (r\'es. $5-3=2$), mais $a \to d$ (r\'es. $3-3=0$) bloqu\'e. $s \to b_{in}$ (r\'es. $4-2=2$) mais $b_{in} \to b_{out}$ (r\'es. $2-2=0$) bloqu\'e. $s \to c$ (r\'es. $3-3=0$) bloqu\'e.

$\Rightarrow$ \textbf{Nouveau flot max $= 8$} (vs 10 sans contrainte nodale).
\end{correction}

\begin{attention}
\begin{itemize}
    \item Le d\'edoublement fonctionne aussi pour les \textbf{chemins sommet-disjoints} : mettre capacit\'e $= 1$ sur chaque arc $v_{in} \to v_{out}$ pour les n\oe uds interm\'ediaires.
    \item Ne jamais d\'edoubler $s$ ni $t$.
\end{itemize}
\end{attention}


\newpage
%=============================================================
%  EXERCICE 7 : CDS + MPR
%=============================================================
\section{Exercice 7 --- Dominant connexe (CDS) et MPR}

\begin{enonce}
Soit le graphe non orient\'e $G$ (r\'eseau sans fil) :

\begin{center}
\begin{tikzpicture}[
    every node/.style={circle, draw, minimum size=9mm, font=\small\bfseries},
    edge/.style={thick}
]
    \node (1) at (0,1) {1};
    \node (2) at (2,2) {2};
    \node (3) at (2,0) {3};
    \node (4) at (4,2) {4};
    \node (5) at (4,0) {5};
    \node (6) at (6,1) {6};
    \node (7) at (6,-1) {7};
    \node (8) at (8,0) {8};

    \draw[edge] (1)--(2); \draw[edge] (1)--(3);
    \draw[edge] (2)--(3); \draw[edge] (2)--(4);
    \draw[edge] (3)--(5); \draw[edge] (4)--(5);
    \draw[edge] (4)--(6); \draw[edge] (5)--(6);
    \draw[edge] (5)--(7); \draw[edge] (6)--(8);
    \draw[edge] (7)--(8);
\end{tikzpicture}
\end{center}

\begin{enumerate}
    \item Appliquer l'algorithme glouton pour construire un dominant connexe (CDS).
    \item Calculer MPR(5) : l'ensemble MPR du n\oe ud 5.
    \item Simuler la diffusion depuis le n\oe ud 1 via le CDS.
\end{enumerate}
\end{enonce}

\begin{methode}
\textbf{CDS glouton :} D\'emarrer avec le n\oe ud de degr\'e max. Ajouter it\'erativement un \textbf{voisin de D} qui couvre le plus de non-domin\'es.

\textbf{MPR de $u$ :} S\'electionner un sous-ensemble minimal de $N(u)$ couvrant tout $N_2(u)$.
\end{methode}

\begin{correction}{title=Correction}

\textbf{Q1. Construction du CDS :}

Degr\'es : $d(1)=2$, $d(2)=3$, $d(3)=3$, $d(4)=3$, $\mathbf{d(5)=4}$, $d(6)=3$, $d(7)=2$, $d(8)=2$.

\begin{etape}{title=\'Etape 1 : Degr\'e max $= 5$ ($d(5)=4$)}
$D = \{5\}$, Domin\'e $= N[5] = \{3, 4, 5, 6, 7\}$.

Non domin\'es : $\{1, 2, 8\}$.
\end{etape}

\begin{etape}{title=\'Etape 2 : Voisins de $D$ non dans $D$ : $\{3, 4, 6, 7\}$}
\begin{itemize}[nosep]
    \item $3$ : couvre $1$ ($\in$ non-domin\'es), $2$ ($\in$ non-domin\'es) $\Rightarrow$ \textbf{2 nouveaux}
    \item $4$ : couvre $2$ $\Rightarrow$ 1 nouveau
    \item $6$ : couvre $8$ $\Rightarrow$ 1 nouveau
    \item $7$ : couvre $8$ $\Rightarrow$ 1 nouveau
\end{itemize}
\textbf{Max :} $3$ (couvre 2 nouveaux). $D = \{3, 5\}$, Domin\'e $= \{1,2,3,4,5,6,7\}$. Non domin\'es : $\{8\}$.
\end{etape}

\begin{etape}{title=\'Etape 3 : Voisins de $D$ non dans $D$ : $\{1, 2, 4, 6, 7\}$}
\begin{itemize}[nosep]
    \item $6$ : couvre $8$ $\Rightarrow$ 1 nouveau
    \item $7$ : couvre $8$ $\Rightarrow$ 1 nouveau
    \item Les autres ne couvrent aucun nouveau.
\end{itemize}
Choisir $6$ (ou $7$). Prenons $6$. $D = \{3, 5, 6\}$. Domin\'e $= V$. \checkmark
\end{etape}

\textbf{CDS} $= \{3, 5, 6\}$. V\'erification :
\begin{itemize}[nosep]
    \item Dominant : $1 \in N(3)$, $2 \in N(3)$, $4 \in N(5)$, $7 \in N(5)$, $8 \in N(6)$ \checkmark
    \item Connexe : $3$-$5$-$6$ forment un chemin dans $G$ \checkmark
\end{itemize}

\medskip
\textbf{Q2. Calcul de MPR(5) :}

$N(5) = \{3, 4, 6, 7\}$

$N_2(5) = $ noeuds \`a distance 2 de 5, pas dans $N[5]$ :
\begin{itemize}[nosep]
    \item Via $3$ : $N(3) = \{1, 2, 5\}$ $\Rightarrow$ $1, 2$ sont dans $N_2(5)$
    \item Via $4$ : $N(4) = \{2, 5, 6\}$ $\Rightarrow$ $2$ d\'ej\`a compt\'e
    \item Via $6$ : $N(6) = \{4, 5, 8\}$ $\Rightarrow$ $8 \in N_2(5)$
    \item Via $7$ : $N(7) = \{5, 8\}$ $\Rightarrow$ $8$ d\'ej\`a compt\'e
\end{itemize}

$N_2(5) = \{1, 2, 8\}$.

\'Etape obligatoire :
\begin{itemize}[nosep]
    \item $1$ : seul voisin commun avec $5$ = $3$ $\Rightarrow$ $3$ obligatoire dans MPR(5)
    \item $2$ : voisins communs avec $5$ = $3$ et $4$ $\Rightarrow$ pas obligatoire
    \item $8$ : voisins communs avec $5$ = $6$ et $7$ $\Rightarrow$ pas obligatoire
\end{itemize}

$\text{MPR}(5) = \{3\}$. Couvert $= \{1, 2\}$ (car $N(3) \cap N_2(5) = \{1, 2\}$).

Non couvert : $\{8\}$. Choisir entre $6$ ($8 \in N(6)$) et $7$ ($8 \in N(7)$). Prenons $6$.

$$\boxed{\text{MPR}(5) = \{3, 6\}}$$

\medskip
\textbf{Q3. Diffusion depuis $1$ via CDS $= \{3, 5, 6\}$ :}

\begin{enumerate}[nosep]
    \item $1$ envoie le message. $1 \notin D$, donc $1$ ne retransmet que s'il est la source (oui).
    \item $2$ et $3$ re\c{c}oivent. $3 \in D$ : \textbf{retransmet}. $2 \notin D$ : ne retransmet pas.
    \item $5$ re\c{c}oit de $3$. $5 \in D$ : \textbf{retransmet}. ($1$ et $2$ ignorent car d\'ej\`a re\c{c}u.)
    \item $4, 6, 7$ re\c{c}oivent de $5$. $6 \in D$ : \textbf{retransmet}. $4, 7 \notin D$ : ne retransmettent pas.
    \item $8$ re\c{c}oit de $6$. $8 \notin D$ : ne retransmet pas.
\end{enumerate}

\textbf{Retransmissions :} $1$ (source) $+ 3 + 5 + 6 = $ \textbf{4 transmissions} pour couvrir 8 n\oe uds.

Flooding aurait donn\'e $8$ transmissions (chaque n\oe ud retransmet une fois).
\end{correction}

\begin{attention}
\begin{itemize}
    \item Pour le CDS : toujours choisir un \textbf{voisin de D} (pas n'importe quel sommet) pour garantir la connexit\'e.
    \item Pour les MPR : ne pas oublier l'\textbf{\'etape obligatoire} (n\oe uds de $N_2$ n'ayant qu'un seul voisin dans $N$).
    \item Le CDS n'est pas unique : diff\'erents choix initiaux donnent diff\'erents CDS.
\end{itemize}
\end{attention}


\newpage
%=============================================================
%  EXERCICE 8 : BELLMAN-FORD
%=============================================================
\section{Exercice 8 --- Bellman-Ford (poids n\'egatifs, graphe g\'en\'eral)}

\begin{enonce}
Soit le graphe orient\'e $G = (V, E)$ avec $V = \{s, 1, 2, 3, 4, 5\}$ et les arcs suivants :

\begin{center}
\begin{tikzpicture}[
    node distance=2.5cm,
    every node/.style={circle, draw, minimum size=9mm, font=\small\bfseries},
    arr/.style={-{Stealth}, thick}
]
    \node[fill=titrebleu!20] (s) at (0,0) {$s$};
    \node[fill=defvert!20] (1) at (2.5,1.5) {$1$};
    \node[fill=defvert!20] (2) at (2.5,-1.5) {$2$};
    \node[fill=exempleorange!20] (3) at (5.5,1.5) {$3$};
    \node[fill=exempleorange!20] (4) at (5.5,-1.5) {$4$};
    \node[fill=importantrouge!20] (5) at (8,0) {$5$};

    \draw[arr] (s) -- node[above left,draw=none,fill=none]{\small $6$} (1);
    \draw[arr] (s) -- node[below left,draw=none,fill=none]{\small $4$} (2);
    \draw[arr] (1) -- node[above,draw=none,fill=none]{\small $5$} (3);
    \draw[arr] (1) -- node[left,draw=none,fill=none]{\small $-2$} (2);
    \draw[arr] (2) -- node[below,draw=none,fill=none]{\small $3$} (4);
    \draw[arr] (3) -- node[above right,draw=none,fill=none]{\small $2$} (5);
    \draw[arr] (4) -- node[below right,draw=none,fill=none]{\small $1$} (5);
    \draw[arr] (3) -- node[right,draw=none,fill=none]{\small $-4$} (4);
    \draw[arr] (4) -- node[below,draw=none,fill=none]{\small $-1$} (2);
\end{tikzpicture}
\end{center}

\begin{enumerate}
    \item Peut-on appliquer Dijkstra ? Pourquoi ?
    \item Ce graphe est-il un DAG ? Peut-on utiliser la relaxation par niveaux ?
    \item Appliquer l'algorithme de \textbf{Bellman-Ford} depuis $s$. Montrer toutes les it\'erations.
    \item Y a-t-il un circuit de poids n\'egatif (circuit absorbant) ?
    \item Donner l'arborescence des plus courts chemins.
\end{enumerate}
\end{enonce}

\begin{methode}
\textbf{Quand utiliser Bellman-Ford ?}
\begin{itemize}[nosep]
    \item Poids n\'egatifs ET le graphe n'est \textbf{pas un DAG} (contient des circuits)
    \item Complexit\'e : $O(|V| \cdot |E|)$
    \item Effectuer $|V|-1$ it\'erations de relaxation de \textbf{tous} les arcs
    \item Une $|V|$-i\`eme it\'eration permet de \textbf{d\'etecter les circuits absorbants} : si un $d[v]$ diminue encore, il y a un circuit de poids n\'egatif
\end{itemize}

\textbf{Proc\'edure :}
\begin{enumerate}[nosep]
    \item Initialiser $d[s] = 0$, $d[v] = +\infty$ pour $v \neq s$, $\pi[v] = \text{nil}$
    \item Pour $k = 1$ \`a $|V|-1$ : pour \textbf{chaque arc} $(u,v)$ de poids $w(u,v)$ : si $d[u] + w(u,v) < d[v]$, alors $d[v] \leftarrow d[u] + w(u,v)$ et $\pi[v] \leftarrow u$
    \item V\'erification : parcourir tous les arcs une derni\`ere fois. Si une am\'elioration est possible $\Rightarrow$ circuit absorbant d\'etect\'e
\end{enumerate}
\end{methode}

\begin{correction}
\textbf{Q1. Dijkstra ?} \textbf{Non.} L'arc $(1,2)$ a un poids $-2$ et l'arc $(3,4)$ a un poids $-4$. Dijkstra \textbf{exige des poids $\geq 0$}.

\medskip
\textbf{Q2. DAG ?} \textbf{Non.} Le graphe contient le circuit $2 \to 4 \to 2$ (via $2 \xrightarrow{3} 4 \xrightarrow{-1} 2$). Poids du circuit $= 3 + (-1) = 2 > 0$, donc ce n'est pas absorbant, mais la pr\'esence d'un circuit emp\^eche d'utiliser la relaxation par niveaux.

$\Rightarrow$ \textbf{Bellman-Ford} est le seul choix.

\medskip
\textbf{Q3. Ex\'ecution de Bellman-Ford.}

Arcs (ordre fix\'e) : $(s,1)_6$, $(s,2)_4$, $(1,3)_5$, $(1,2)_{-2}$, $(2,4)_3$, $(3,5)_2$, $(3,4)_{-4}$, $(4,5)_1$, $(4,2)_{-1}$.

\medskip
\textbf{Initialisation :}

\begin{center}
\renewcommand{\arraystretch}{1.2}
\begin{tabular}{|c|c|c|c|c|c|c|}
\hline
& $s$ & $1$ & $2$ & $3$ & $4$ & $5$ \\\hline
Init & $0$ & $+\infty$ & $+\infty$ & $+\infty$ & $+\infty$ & $+\infty$ \\\hline
\end{tabular}
\end{center}

\textbf{It\'eration 1} ($k=1$) --- relaxer tous les arcs :
\begin{itemize}[nosep]
    \item $(s,1)$ : $d[1] = \min(\infty, 0+6) = \mathbf{6}$, $\pi[1]=s$
    \item $(s,2)$ : $d[2] = \min(\infty, 0+4) = \mathbf{4}$, $\pi[2]=s$
    \item $(1,3)$ : $d[3] = \min(\infty, 6+5) = \mathbf{11}$, $\pi[3]=1$
    \item $(1,2)$ : $d[2] = \min(4, 6+(-2)) = \mathbf{4}$ (inchang\'e)
    \item $(2,4)$ : $d[4] = \min(\infty, 4+3) = \mathbf{7}$, $\pi[4]=2$
    \item $(3,5)$ : $d[5] = \min(\infty, 11+2) = \mathbf{13}$, $\pi[5]=3$
    \item $(3,4)$ : $d[4] = \min(7, 11+(-4)) = \mathbf{7}$ (inchang\'e)
    \item $(4,5)$ : $d[5] = \min(13, 7+1) = \mathbf{8}$, $\pi[5]=4$
    \item $(4,2)$ : $d[2] = \min(4, 7+(-1)) = \mathbf{4}$ (inchang\'e)
\end{itemize}

\begin{center}
\begin{tabular}{|c|c|c|c|c|c|c|}
\hline
& $s$ & $1$ & $2$ & $3$ & $4$ & $5$ \\\hline
$k=1$ & $0$ & $6$ & $4$ & $11$ & $7$ & $8$ \\\hline
\end{tabular}
\end{center}

\textbf{It\'eration 2} ($k=2$) --- relaxer tous les arcs :
\begin{itemize}[nosep]
    \item $(s,1)$ : $d[1] = \min(6, 0+6) = 6$ (inchang\'e)
    \item $(s,2)$ : $d[2] = \min(4, 0+4) = 4$ (inchang\'e)
    \item $(1,3)$ : $d[3] = \min(11, 6+5) = 11$ (inchang\'e)
    \item $(1,2)$ : $d[2] = \min(4, 6-2) = \mathbf{4}$ (inchang\'e)
    \item $(2,4)$ : $d[4] = \min(7, 4+3) = \mathbf{7}$ (inchang\'e)
    \item $(3,5)$ : $d[5] = \min(8, 11+2) = 8$ (inchang\'e)
    \item $(3,4)$ : $d[4] = \min(7, 11-4) = \mathbf{7}$ (inchang\'e)
    \item $(4,5)$ : $d[5] = \min(8, 7+1) = \mathbf{8}$ (inchang\'e)
    \item $(4,2)$ : $d[2] = \min(4, 7-1) = \mathbf{4}$ (inchang\'e)
\end{itemize}

\begin{center}
\begin{tabular}{|c|c|c|c|c|c|c|}
\hline
& $s$ & $1$ & $2$ & $3$ & $4$ & $5$ \\\hline
$k=2$ & $0$ & $6$ & $4$ & $11$ & $7$ & $8$ \\\hline
\end{tabular}
\end{center}

\textbf{It\'erations 3, 4, 5} ($k=3$ \`a $|V|-1=5$) : aucune modification $\Rightarrow$ convergence atteinte d\`es $k=1$.

\medskip
\textbf{Q4. D\'etection de circuit absorbant.}

It\'eration $|V| = 6$ (v\'erification) : on relaxe tous les arcs une derni\`ere fois.

\textbf{Aucune am\'elioration} $\Rightarrow$ \textcolor{defvert}{\textbf{pas de circuit absorbant}}.

(On peut v\'erifier : le seul circuit est $2 \to 4 \to 2$ de poids $3 + (-1) = 2 > 0$.)

\medskip
\textbf{Q5. Arborescence des PCC :}

$$\pi[1]=s,\quad \pi[2]=s,\quad \pi[3]=1,\quad \pi[4]=2,\quad \pi[5]=4$$

\begin{center}
\begin{tikzpicture}[
    node distance=2cm,
    every node/.style={circle, draw, minimum size=9mm, font=\small\bfseries},
    arr/.style={-{Stealth}, thick, color=defvert}
]
    \node[fill=titrebleu!20] (s) at (0,0) {$s$};
    \node[fill=defvert!20] (1) at (2.5,1.2) {$1$};
    \node[fill=defvert!20] (2) at (2.5,-1.2) {$2$};
    \node[fill=exempleorange!20] (3) at (5,1.2) {$3$};
    \node[fill=exempleorange!20] (4) at (5,-1.2) {$4$};
    \node[fill=importantrouge!20] (5) at (7.5,0) {$5$};

    \draw[arr] (s) -- node[above left,draw=none,fill=none]{\small $d=6$} (1);
    \draw[arr] (s) -- node[below left,draw=none,fill=none]{\small $d=4$} (2);
    \draw[arr] (1) -- node[above,draw=none,fill=none]{\small $d=11$} (3);
    \draw[arr] (2) -- node[below,draw=none,fill=none]{\small $d=7$} (4);
    \draw[arr] (4) -- node[below right,draw=none,fill=none]{\small $d=8$} (5);
\end{tikzpicture}
\end{center}

\textbf{Chemins optimaux :}
\begin{align*}
s &\to 1 \quad (d=6) \\
s &\to 2 \quad (d=4) \\
s &\to 1 \to 3 \quad (d=11) \\
s &\to 2 \to 4 \quad (d=7) \\
s &\to 2 \to 4 \to 5 \quad (d=8)
\end{align*}
\end{correction}

\begin{attention}
\begin{itemize}
    \item \textbf{Bellman-Ford vs Dijkstra :} Bellman-Ford est le seul choix quand il y a des poids n\'egatifs ET que le graphe n'est pas un DAG.
    \item \textbf{Nombre d'it\'erations :} exactement $|V|-1$ it\'erations (m\^eme si convergence avant). La $|V|$-i\`eme it\'eration sert \textbf{uniquement} \`a d\'etecter les circuits absorbants.
    \item \textbf{Ordre des arcs :} l'ordre de parcours des arcs dans chaque it\'eration ne change pas le r\'esultat final, mais peut acc\'el\'erer la convergence.
    \item \textbf{Circuit de poids n\'egatif $\neq$ arc de poids n\'egatif :} un arc n\'egatif ne pose pas de probl\`eme en soi. C'est un \textbf{circuit} dont la somme des poids est $< 0$ qui rend le probl\`eme insoluble.
\end{itemize}
\end{attention}


\newpage
%=============================================================
%  EXERCICE 9 : FORD-FULKERSON AVEC ARCS RETOUR
%=============================================================
\section{Exercice 9 --- Flot max avec chemins augmentants et arcs retour}

\begin{enonce}
Soit le r\'eseau $G$ suivant avec source $s$ et puits $t$. Les valeurs sur les arcs indiquent la capacit\'e.

\begin{center}
\begin{tikzpicture}[
    node distance=2.5cm,
    every node/.style={circle, draw, minimum size=9mm, font=\small\bfseries},
    arr/.style={-{Stealth}, thick}
]
    \node[fill=titrebleu!20] (s) at (0,0) {$s$};
    \node[fill=defvert!20] (a) at (2.5,1.5) {$a$};
    \node[fill=defvert!20] (b) at (2.5,-1.5) {$b$};
    \node[fill=exempleorange!20] (c) at (5.5,1.5) {$c$};
    \node[fill=exempleorange!20] (d) at (5.5,-1.5) {$d$};
    \node[fill=importantrouge!20] (t) at (8,0) {$t$};

    \draw[arr] (s) -- node[above left,draw=none,fill=none]{\small $10$} (a);
    \draw[arr] (s) -- node[below left,draw=none,fill=none]{\small $8$} (b);
    \draw[arr] (a) -- node[above,draw=none,fill=none]{\small $6$} (c);
    \draw[arr] (a) -- node[left,draw=none,fill=none]{\small $5$} (d);
    \draw[arr] (b) -- node[below,draw=none,fill=none]{\small $7$} (d);
    \draw[arr] (b) -- node[right,draw=none,fill=none]{\small $3$} (a);
    \draw[arr] (c) -- node[above right,draw=none,fill=none]{\small $9$} (t);
    \draw[arr] (d) -- node[below right,draw=none,fill=none]{\small $10$} (t);
    \draw[arr] (c) -- node[right,draw=none,fill=none]{\small $4$} (d);
\end{tikzpicture}
\end{center}

\begin{enumerate}
    \item Trouver un flot initial en identifiant des chemins augmentants \'evidents.
    \item Construire le graphe r\'esiduel apr\`es ce flot initial.
    \item Chercher un nouveau chemin augmentant (qui utilise un \textbf{arc retour} dans le graphe r\'esiduel). Augmenter le flot.
    \item R\'ep\'eter jusqu'\`a obtenir le flot maximum. Donner $f^*$.
    \item Identifier une coupe minimum et v\'erifier le th\'eor\`eme flot max / coupe min.
\end{enumerate}
\end{enonce}

\begin{methode}
\textbf{Ford-Fulkerson --- graphe r\'esiduel :}
\begin{itemize}[nosep]
    \item Pour chaque arc $(u,v)$ de capacit\'e $c$ portant un flot $f$ :
    \begin{itemize}[nosep]
        \item \textbf{Arc avant} $(u \to v)$ : capacit\'e r\'esiduelle $= c - f$ (si $> 0$)
        \item \textbf{Arc retour} $(v \to u)$ : capacit\'e r\'esiduelle $= f$ (si $> 0$)
    \end{itemize}
    \item Un \textbf{chemin augmentant} utilisant un arc retour $(v,u)$ signifie qu'on \textbf{diminue} le flot sur l'arc original $(u,v)$ pour le rediriger ailleurs.
    \item La valeur d'augmentation = $\min$ des capacit\'es r\'esiduelles sur le chemin.
\end{itemize}
\end{methode}

\begin{correction}
\textbf{Q1. Flot initial --- chemins augmentants \'evidents :}

\begin{itemize}[nosep]
    \item $P_1 : s \to a \to c \to t$ : goulot $= \min(10, 6, 9) = 6$. Envoyer $f=6$.
    \item $P_2 : s \to b \to d \to t$ : goulot $= \min(8, 7, 10) = 7$. Envoyer $f=7$.
\end{itemize}

\'Etat du flot apr\`es $P_1$ et $P_2$ (notation $f/c$) :

\begin{center}
\begin{tabular}{|c|c|c|c|c|c|c|c|c|c|}
\hline
Arc & $s\!\to\!a$ & $s\!\to\!b$ & $a\!\to\!c$ & $a\!\to\!d$ & $b\!\to\!d$ & $b\!\to\!a$ & $c\!\to\!t$ & $d\!\to\!t$ & $c\!\to\!d$ \\\hline
$f/c$ & $6/10$ & $7/8$ & $6/6$ & $0/5$ & $7/7$ & $0/3$ & $6/9$ & $7/10$ & $0/4$ \\\hline
\end{tabular}
\end{center}

Flot actuel : $|f| = 6 + 7 = 13$.

\medskip
\textbf{Q2. Graphe r\'esiduel $G^r$ :}

\begin{center}
\begin{tikzpicture}[
    node distance=2.5cm,
    every node/.style={circle, draw, minimum size=9mm, font=\small\bfseries},
    arr/.style={-{Stealth}, thick, color=titrebleu},
    ret/.style={-{Stealth}, thick, color=importantrouge, dashed}
]
    \node[fill=titrebleu!20] (s) at (0,0) {$s$};
    \node[fill=defvert!20] (a) at (2.5,1.5) {$a$};
    \node[fill=defvert!20] (b) at (2.5,-1.5) {$b$};
    \node[fill=exempleorange!20] (c) at (5.5,1.5) {$c$};
    \node[fill=exempleorange!20] (d) at (5.5,-1.5) {$d$};
    \node[fill=importantrouge!20] (t) at (8,0) {$t$};

    % Arcs avant (capacit\'e r\'esiduelle > 0)
    \draw[arr] (s) to[bend left=10] node[above left,draw=none,fill=none]{\small $4$} (a);
    \draw[arr] (s) to[bend right=10] node[below left,draw=none,fill=none]{\small $1$} (b);
    \draw[arr] (a) -- node[left,draw=none,fill=none]{\small $5$} (d);
    \draw[arr] (b) -- node[right,draw=none,fill=none]{\small $3$} (a);
    \draw[arr] (c) to[bend left=10] node[above right,draw=none,fill=none]{\small $3$} (t);
    \draw[arr] (d) to[bend right=10] node[below right,draw=none,fill=none]{\small $3$} (t);
    \draw[arr] (c) -- node[right,draw=none,fill=none]{\small $4$} (d);

    % Arcs retour
    \draw[ret] (a) to[bend left=10] node[below right,draw=none,fill=none]{\small $6$} (s);
    \draw[ret] (b) to[bend right=10] node[above right,draw=none,fill=none]{\small $7$} (s);
    \draw[ret] (c) -- node[below,draw=none,fill=none]{\small $6$} (a);
    \draw[ret] (d) to[bend left=10] node[above,draw=none,fill=none]{\small $7$} (b);
    \draw[ret] (t) to[bend left=10] node[below left,draw=none,fill=none]{\small $6$} (c);
    \draw[ret] (t) to[bend right=10] node[above left,draw=none,fill=none]{\small $7$} (d);
\end{tikzpicture}
\end{center}

{\small Bleu = arcs avant (capacit\'e r\'esiduelle $c-f$). Rouge pointill\'e = arcs retour (capacit\'e $f$).}

\medskip
\textbf{Q3. Chemin augmentant avec arc retour :}

$P_3 : s \to b \to a \to d \to t$ dans $G^r$.

Capacit\'es r\'esiduelles : $\min(1, 3, 5, 3) = 1$. Augmenter de $1$.

Mise \`a jour :
\begin{itemize}[nosep]
    \item $s \to b$ : $f = 7+1 = 8$ (satur\'e)
    \item $b \to a$ : $f = 0+1 = 1$
    \item $a \to d$ : $f = 0+1 = 1$
    \item $d \to t$ : $f = 7+1 = 8$
\end{itemize}

Flot : $|f| = 13 + 1 = 14$.

\medskip
\textbf{Nouveau graphe r\'esiduel :} Cherchons un nouveau chemin augmentant.

$P_4 : s \to a \to c \to d \to t$ dans $G^r$.

Capacit\'es r\'esiduelles : $s \to a : 10-6 = 4 - 0 = 3$ (puisque $f(s,a)=6+0=6$, non, recalculons).

Apr\`es $P_3$, le flot est :

\begin{center}
\begin{tabular}{|c|c|c|c|c|c|c|c|c|c|}
\hline
Arc & $s\!\to\!a$ & $s\!\to\!b$ & $a\!\to\!c$ & $a\!\to\!d$ & $b\!\to\!d$ & $b\!\to\!a$ & $c\!\to\!t$ & $d\!\to\!t$ & $c\!\to\!d$ \\\hline
$f/c$ & $6/10$ & $8/8$ & $6/6$ & $1/5$ & $7/7$ & $1/3$ & $6/9$ & $8/10$ & $0/4$ \\\hline
\end{tabular}
\end{center}

$P_4 : s \to a \to d \to t$ : $\min(4, 4, 2) = 2$. Augmenter de $2$.

Apr\`es $P_4$ :

\begin{center}
\begin{tabular}{|c|c|c|c|c|c|c|c|c|c|}
\hline
Arc & $s\!\to\!a$ & $s\!\to\!b$ & $a\!\to\!c$ & $a\!\to\!d$ & $b\!\to\!d$ & $b\!\to\!a$ & $c\!\to\!t$ & $d\!\to\!t$ & $c\!\to\!d$ \\\hline
$f/c$ & $8/10$ & $8/8$ & $6/6$ & $3/5$ & $7/7$ & $1/3$ & $6/9$ & $10/10$ & $0/4$ \\\hline
\end{tabular}
\end{center}

Flot : $|f| = 14 + 2 = 16$.

\medskip
Nouveau $G^r$ : $d \to t$ satur\'e, $s \to b$ satur\'e. Cherchons un chemin de $s$ \`a $t$ :

$P_5 : s \to a \to c \to t$ ? Non, $a \to c$ satur\'e. Mais arc retour $c \to a$ existe.

$s \to a$ (cap. 2) ? Puis $a \to d$ (cap. 2) ? Puis $d \to t$ (cap. 0) : bloqu\'e.

$s \to a$ (cap. 2) $\to$ $a \to c$? Non (cap. $0$).

Plus de chemin $s$-$t$ dans $G^r$.

\medskip
\textbf{Q4. Flot maximum :}
$$\boxed{f^* = 16}$$

V\'erification : $f^* = f(s,a) + f(s,b) = 8 + 8 = 16$ (sortant de $s$) et $f(c,t) + f(d,t) = 6 + 10 = 16$ (entrant dans $t$). $\checkmark$

\medskip
\textbf{Q5. Coupe minimum :}

Depuis $s$ dans le graphe r\'esiduel final, les sommets accessibles sont $S = \{s, a\}$ (car $s \to a$ a encore cap. r\'esiduelle $= 2$, mais depuis $a$ : $a \to c$ satur\'e, $a \to d$ cap. $2$ $\to d \to t$ satur\'e).

Recalculons : depuis $s$, on peut atteindre $a$ (cap. 2). Depuis $a$ : $a \to d$ (cap. 2) $\to$ on atteint $d$. Depuis $d$ : $d \to t$ satur\'e (cap. 0), arc retour $d \to b$ (cap. 7), arc retour $d \to a$ (cap. 1). Depuis $b$ : arc retour $b \to s$ (cap. 8), $b \to a$ (cap. 2). Donc $S = \{s, a, b, d\}$, $\bar{S} = \{c, t\}$.

Coupe $(S, \bar{S})$ : arcs de $S$ vers $\bar{S}$ :
\begin{itemize}[nosep]
    \item $(a, c)$ : capacit\'e $6$
    \item $(d, t)$ : capacit\'e $10$
\end{itemize}

$$\text{Capacit\'e de la coupe} = 6 + 10 = 16 = f^*\ \checkmark$$

\textbf{Th\'eor\`eme flot max / coupe min v\'erifi\'e.}
\end{correction}

\begin{attention}
\begin{itemize}
    \item Les \textbf{arcs retour} dans le graphe r\'esiduel permettent de \og d\'efaire \fg{} un flot mal plac\'e pour le rediriger ailleurs. C'est souvent n\'ecessaire pour atteindre l'optimum.
    \item La \textbf{coupe minimum} se trouve en faisant un BFS/DFS depuis $s$ dans le graphe r\'esiduel final : $S$ = sommets accessibles, $\bar{S}$ = le reste.
    \item Notation examen : $a/b$ signifie flot $a$ sur capacit\'e $b$. Capacit\'e r\'esiduelle avant $= b-a$, retour $= a$.
\end{itemize}
\end{attention}


\newpage
%=============================================================
%  EXERCICE 10 : CHEMINS SOMMETS-DISJOINTS (DEDOUBLEMENT)
%=============================================================
\section{Exercice 10 --- Chemins sommets-disjoints (d\'edoublement de n\oe uds)}

\begin{enonce}
Soit le graphe orient\'e $G = (V, E)$ suivant avec $V = \{s, a, b, c, d, t\}$ :

\begin{center}
\begin{tikzpicture}[
    node distance=2.5cm,
    every node/.style={circle, draw, minimum size=9mm, font=\small\bfseries},
    arr/.style={-{Stealth}, thick}
]
    \node[fill=titrebleu!20] (s) at (0,0) {$s$};
    \node[fill=defvert!20] (a) at (2.5,1.5) {$a$};
    \node[fill=defvert!20] (b) at (2.5,-1.5) {$b$};
    \node[fill=exempleorange!20] (c) at (5,1.5) {$c$};
    \node[fill=exempleorange!20] (d) at (5,-1.5) {$d$};
    \node[fill=importantrouge!20] (t) at (7.5,0) {$t$};

    \draw[arr] (s) -- (a);
    \draw[arr] (s) -- (b);
    \draw[arr] (a) -- (c);
    \draw[arr] (a) -- (d);
    \draw[arr] (b) -- (c);
    \draw[arr] (b) -- (d);
    \draw[arr] (c) -- (t);
    \draw[arr] (d) -- (t);
\end{tikzpicture}
\end{center}

\begin{enumerate}
    \item Trouver le nombre maximum de chemins \textbf{arcs-disjoints} de $s$ \`a $t$.
    \item Trouver le nombre maximum de chemins \textbf{sommets-disjoints} de $s$ \`a $t$ (les chemins ne partagent aucun sommet interm\'ediaire).
    \item Pour la question 2, appliquer la technique de \textbf{d\'edoublement de n\oe uds} : construire le graphe transform\'e $G'$, puis trouver le flot maximum.
\end{enumerate}
\end{enonce}

\begin{methode}
\textbf{Chemins arcs-disjoints :} flot max avec capacit\'es $= 1$ sur chaque arc (th. de Menger, version arcs).

\textbf{Chemins sommets-disjoints :} d\'edoublement de n\oe uds (th. de Menger, version sommets).
\begin{itemize}[nosep]
    \item Chaque sommet interm\'ediaire $v$ (ni $s$ ni $t$) est remplac\'e par $v_{in}$ et $v_{out}$
    \item Arc interne $(v_{in} \to v_{out})$ de capacit\'e $1$ (un seul chemin peut traverser $v$)
    \item Tout arc $(u, v)$ du graphe original devient $(u_{out}, v_{in})$ de capacit\'e $1$
    \item $s$ garde un seul n\oe ud $s_{out}$, $t$ garde un seul n\oe ud $t_{in}$
    \item Le flot max dans $G'$ donne le nombre max de chemins sommets-disjoints dans $G$
\end{itemize}
\end{methode}

\begin{correction}
\textbf{Q1. Chemins arcs-disjoints :}

Capacit\'es $= 1$ sur tous les arcs. Cherchons des chemins $s$-$t$ arcs-disjoints :
\begin{itemize}[nosep]
    \item $P_1 : s \to a \to c \to t$ \quad (arcs $sa$, $ac$, $ct$)
    \item $P_2 : s \to b \to d \to t$ \quad (arcs $sb$, $bd$, $dt$)
\end{itemize}

Aucun arc commun $\checkmark$. Peut-on faire 3 ? Non : $t$ n'a que 2 arcs entrants ($ct$ et $dt$), donc au plus 2 chemins.

$$\boxed{\text{Chemins arcs-disjoints max} = 2}$$

\medskip
\textbf{Q2. Chemins sommets-disjoints :}

Les chemins $P_1$ et $P_2$ ci-dessus ne partagent aucun sommet interm\'ediaire ($a, c$ pour $P_1$ et $b, d$ pour $P_2$) $\Rightarrow$ ils sont aussi sommets-disjoints.

$$\boxed{\text{Chemins sommets-disjoints max} = 2}$$

\medskip
\textbf{Q3. V\'erification par d\'edoublement :}

Construction de $G'$ : chaque sommet interm\'ediaire $v \in \{a, b, c, d\}$ est remplac\'e par $(v_{in}, v_{out})$ avec un arc de capacit\'e $1$.

\begin{center}
\begin{tikzpicture}[
    node distance=1.8cm,
    every node/.style={circle, draw, minimum size=8mm, font=\scriptsize\bfseries},
    arr/.style={-{Stealth}, thick},
    internal/.style={-{Stealth}, thick, color=importantrouge}
]
    \node[fill=titrebleu!20] (s) at (0,0) {$s$};

    \node[fill=defvert!20] (ain) at (2,1.8) {$a_i$};
    \node[fill=defvert!20] (aout) at (3.5,1.8) {$a_o$};

    \node[fill=defvert!20] (bin) at (2,-1.8) {$b_i$};
    \node[fill=defvert!20] (bout) at (3.5,-1.8) {$b_o$};

    \node[fill=exempleorange!20] (cin) at (5.5,1.8) {$c_i$};
    \node[fill=exempleorange!20] (cout) at (7,1.8) {$c_o$};

    \node[fill=exempleorange!20] (din) at (5.5,-1.8) {$d_i$};
    \node[fill=exempleorange!20] (dout) at (7,-1.8) {$d_o$};

    \node[fill=importantrouge!20] (t) at (9,0) {$t$};

    % Arcs internes (capacit\'e 1)
    \draw[internal] (ain) -- node[above,draw=none,fill=none]{\tiny $1$} (aout);
    \draw[internal] (bin) -- node[above,draw=none,fill=none]{\tiny $1$} (bout);
    \draw[internal] (cin) -- node[above,draw=none,fill=none]{\tiny $1$} (cout);
    \draw[internal] (din) -- node[above,draw=none,fill=none]{\tiny $1$} (dout);

    % Arcs du graphe original
    \draw[arr] (s) -- node[above,draw=none,fill=none]{\tiny $1$} (ain);
    \draw[arr] (s) -- node[below,draw=none,fill=none]{\tiny $1$} (bin);
    \draw[arr] (aout) -- node[above,draw=none,fill=none]{\tiny $1$} (cin);
    \draw[arr] (aout) -- node[left,draw=none,fill=none]{\tiny $1$} (din);
    \draw[arr] (bout) -- node[right,draw=none,fill=none]{\tiny $1$} (cin);
    \draw[arr] (bout) -- node[below,draw=none,fill=none]{\tiny $1$} (din);
    \draw[arr] (cout) -- node[above,draw=none,fill=none]{\tiny $1$} (t);
    \draw[arr] (dout) -- node[below,draw=none,fill=none]{\tiny $1$} (t);
\end{tikzpicture}
\end{center}

{\small Arcs rouges = arcs internes $(v_{in} \to v_{out})$ de capacit\'e 1. Arcs noirs = arcs originaux de capacit\'e 1.}

\medskip
Flot max dans $G'$ :
\begin{itemize}[nosep]
    \item $P_1' : s \to a_{in} \to a_{out} \to c_{in} \to c_{out} \to t$ (flot $= 1$)
    \item $P_2' : s \to b_{in} \to b_{out} \to d_{in} \to d_{out} \to t$ (flot $= 1$)
\end{itemize}

$|f^*| = 2$. Les arcs internes de $a$ et $b$ sont satur\'es. Un 3\textsuperscript{e} chemin devrait passer par $a$ ou $b$ (seuls successeurs de $s$), mais $a_{in} \to a_{out}$ et $b_{in} \to b_{out}$ sont satur\'es.

$$\boxed{\text{Nombre max de chemins sommets-disjoints} = 2}$$

Confirmation du th\'eor\`eme de Menger (version sommets) : le nombre maximum de chemins sommets-disjoints de $s$ \`a $t$ = taille minimum d'un ensemble de sommets (hors $s,t$) dont la suppression d\'econnecte $s$ de $t$. Ici, supprimer $\{a, b\}$ (ou $\{c, d\}$) d\'econnecte $s$ de $t$ : connectivit\'e nodale $\kappa(s,t) = 2$.
\end{correction}

\begin{attention}
\begin{itemize}
    \item \textbf{Ne pas d\'edoubler $s$ et $t$} : ils ne subissent pas la contrainte de capacit\'e nodale.
    \item \textbf{Capacit\'e nodale $k$} : si on veut que chaque sommet supporte au plus $k$ chemins, mettre la capacit\'e de l'arc interne \`a $k$ (au lieu de $1$).
    \item Toujours v\'erifier : chemins sommets-disjoints $\leq$ chemins arcs-disjoints (car disjoints sur les sommets $\Rightarrow$ disjoints sur les arcs, mais pas l'inverse).
\end{itemize}
\end{attention}


\newpage
%=============================================================
%  EXERCICE 11 : EXERCICE TYPE EXAMEN (COMBINE)
%=============================================================
\section{Exercice 11 --- Exercice type examen (synth\`ese compl\`ete)}

\begin{enonce}
Soit le graphe orient\'e $G = (V, E)$ avec $V = \{s, 1, 2, 3, 4, 5, t\}$ et les arcs pond\'er\'es :

\begin{center}
\begin{tikzpicture}[
    node distance=2.5cm,
    every node/.style={circle, draw, minimum size=9mm, font=\small\bfseries},
    arr/.style={-{Stealth}, thick}
]
    \node[fill=titrebleu!20] (s) at (0,0) {$s$};
    \node[fill=defvert!20] (1) at (2.5,2) {$1$};
    \node[fill=defvert!20] (2) at (2.5,-2) {$2$};
    \node[fill=exempleorange!20] (3) at (5,2) {$3$};
    \node[fill=exempleorange!20] (4) at (5,0) {$4$};
    \node[fill=algoviolet!20] (5) at (5,-2) {$5$};
    \node[fill=importantrouge!20] (t) at (7.5,0) {$t$};

    \draw[arr] (s) -- node[above left,draw=none,fill=none]{\small $3$} (1);
    \draw[arr] (s) -- node[left,draw=none,fill=none]{\small $5$} (4);
    \draw[arr] (s) -- node[below left,draw=none,fill=none]{\small $2$} (2);
    \draw[arr] (1) -- node[above,draw=none,fill=none]{\small $1$} (3);
    \draw[arr] (1) -- node[left,draw=none,fill=none]{\small $4$} (4);
    \draw[arr] (2) -- node[below,draw=none,fill=none]{\small $3$} (5);
    \draw[arr] (2) -- node[right,draw=none,fill=none]{\small $6$} (4);
    \draw[arr] (3) -- node[above right,draw=none,fill=none]{\small $2$} (t);
    \draw[arr] (4) -- node[right,draw=none,fill=none]{\small $1$} (t);
    \draw[arr] (5) -- node[below right,draw=none,fill=none]{\small $4$} (t);
    \draw[arr] (3) -- node[right,draw=none,fill=none]{\small $3$} (4);
\end{tikzpicture}
\end{center}

\textbf{Partie A --- Plus courts chemins :}
\begin{enumerate}
    \item Ce graphe est-il un DAG ? Si oui, donner les niveaux.
    \item Quel algorithme choisir pour calculer les PCC depuis $s$ ? Justifier.
    \item Ex\'ecuter l'algorithme choisi. Donner les distances et l'arborescence.
\end{enumerate}

\textbf{Partie B --- R\'eseau sans fil :}

On consid\`ere maintenant le graphe $G$ comme \textbf{non orient\'e} (les poids deviennent des co\^uts d'\'emission).
\begin{enumerate}[resume]
    \item Trouver un arbre couvrant de poids minimum (MST) par Kruskal.
    \item Calculer la puissance de chaque n\oe ud sur le MST et la puissance totale.
    \item Appliquer l'algorithme BIP depuis $s$ pour construire un arbre de diffusion.
    \item Comparer la puissance totale BIP avec celle du MST.
\end{enumerate}
\end{enonce}

\begin{correction}
%--- PARTIE A ---
\textbf{Partie A --- Plus courts chemins}

\medskip
\textbf{Q1. DAG et niveaux ?}

V\'erifions l'absence de circuit : tous les arcs vont de $s$ vers $\{1,2,4\}$, de $\{1,2\}$ vers $\{3,4,5\}$, et de $\{3,4,5\}$ vers $t$. Aucun arc ne revient en arri\`ere $\Rightarrow$ \textbf{G est un DAG}.

Niveaux (par calcul des pr\'ed\'ecesseurs) :
\begin{itemize}[nosep]
    \item $s$ : aucun pr\'ed\'ecesseur $\Rightarrow$ \textbf{Niveau 0}
    \item $1$ : pr\'edec. $s$ (niv.0) $\Rightarrow$ \textbf{Niveau 1}
    \item $2$ : pr\'edec. $s$ (niv.0) $\Rightarrow$ \textbf{Niveau 1}
    \item $4$ : pr\'edec. $s$(0), $1$(1), $2$(1), $3$(2) $\Rightarrow$ \textbf{Niveau 3}
    \item $3$ : pr\'edec. $1$(1) $\Rightarrow$ \textbf{Niveau 2}
    \item $5$ : pr\'edec. $2$(1) $\Rightarrow$ \textbf{Niveau 2}
    \item $t$ : pr\'edec. $3$(2), $4$(3), $5$(2) $\Rightarrow$ \textbf{Niveau 4}
\end{itemize}

\begin{center}
\begin{tikzpicture}[
    every node/.style={circle, draw, minimum size=8mm, font=\small\bfseries}
]
    \node[fill=titrebleu!20] (s) at (0,0) {$s$};
    \node[fill=defvert!20] (1) at (2,1) {$1$};
    \node[fill=defvert!20] (2) at (2,-1) {$2$};
    \node[fill=exempleorange!20] (3) at (4,1) {$3$};
    \node[fill=exempleorange!20] (5) at (4,-1) {$5$};
    \node[fill=algoviolet!20] (4) at (6,0) {$4$};
    \node[fill=importantrouge!20] (t) at (8,0) {$t$};

    \node[draw=none,fill=none,color=titrebleu] at (0,-2) {\small $L_0$};
    \node[draw=none,fill=none,color=defvert] at (2,-2) {\small $L_1$};
    \node[draw=none,fill=none,color=exempleorange] at (4,-2) {\small $L_2$};
    \node[draw=none,fill=none,color=algoviolet] at (6,-2) {\small $L_3$};
    \node[draw=none,fill=none,color=importantrouge] at (8,-2) {\small $L_4$};
\end{tikzpicture}
\end{center}

\textbf{Q2. Algorithme ?}

Tous les poids sont $\geq 0$ et le graphe est un DAG $\Rightarrow$ on peut utiliser Dijkstra \textbf{ou} la relaxation par niveaux. La \textbf{relaxation par niveaux} est optimale : $O(|V|+|E|)$ contre $O(|V|^2)$ pour Dijkstra.

\textbf{Choix : relaxation topologique (par niveaux).}

\medskip
\textbf{Q3. Ex\'ecution :}

Initialisation : $d[s]=0$, tous les autres $= +\infty$.

\textbf{Niveau 0 --- $s$ ($d=0$) :}
\begin{align*}
d[1] &= \min(\infty, 0+3) = 3,\quad d[4] = \min(\infty, 0+5) = 5,\quad d[2] = \min(\infty, 0+2) = 2
\end{align*}

\textbf{Niveau 1 --- $1$ ($d=3$) puis $2$ ($d=2$) :}

Depuis $1$ : $d[3] = \min(\infty, 3+1) = 4$, $d[4] = \min(5, 3+4) = 5$ (inchang\'e).

Depuis $2$ : $d[5] = \min(\infty, 2+3) = 5$, $d[4] = \min(5, 2+6) = 5$ (inchang\'e).

\textbf{Niveau 2 --- $3$ ($d=4$) puis $5$ ($d=5$) :}

Depuis $3$ : $d[t] = \min(\infty, 4+2) = 6$, $d[4] = \min(5, 4+3) = 5$ (inchang\'e).

Depuis $5$ : $d[t] = \min(6, 5+4) = 6$ (inchang\'e).

\textbf{Niveau 3 --- $4$ ($d=5$) :}

Depuis $4$ : $d[t] = \min(6, 5+1) = \mathbf{6}$ (inchang\'e).

\begin{center}
\renewcommand{\arraystretch}{1.2}
\begin{tabular}{|c|c|c|c|c|c|c|c|}
\hline
& $s$ & $1$ & $2$ & $3$ & $4$ & $5$ & $t$ \\\hline
$d(\cdot)$ & $0$ & $3$ & $2$ & $4$ & $5$ & $5$ & $6$ \\\hline
$\pi(\cdot)$ & --- & $s$ & $s$ & $1$ & $s$ & $2$ & $3$ \\\hline
\end{tabular}
\end{center}

Arborescence : $s \to 1 \to 3 \to t$ ($d=6$), $s \to 2 \to 5$ ($d=5$), $s \to 4$ ($d=5$).

Chemins : $s \xrightarrow{3} 1 \xrightarrow{1} 3 \xrightarrow{2} t$ (distance 6). \'Egalement $s \xrightarrow{5} 4 \xrightarrow{1} t = 6$.

\bigskip
%--- PARTIE B ---
\textbf{Partie B --- R\'eseau sans fil (graphe non orient\'e)}

\medskip
Ar\^etes avec poids : $s$-$1$(3), $s$-$2$(2), $s$-$4$(5), $1$-$3$(1), $1$-$4$(4), $2$-$4$(6), $2$-$5$(3), $3$-$4$(3), $3$-$t$(2), $4$-$t$(1), $5$-$t$(4).

\medskip
\textbf{Q4. MST par Kruskal :}

Ar\^etes tri\'ees par poids croissant :

$1$-$3$(1), $4$-$t$(1), $s$-$2$(2), $3$-$t$(2), $s$-$1$(3), $2$-$5$(3), $3$-$4$(3), $1$-$4$(4), $5$-$t$(4), $s$-$4$(5), $2$-$4$(6)

\begin{enumerate}[nosep]
    \item $1$-$3$ (1) : ajouter $\to \{1,3\}$
    \item $4$-$t$ (1) : ajouter $\to \{4,t\}$
    \item $s$-$2$ (2) : ajouter $\to \{s,2\}$
    \item $3$-$t$ (2) : ajouter $\to \{1,3,4,t\}$ (fusionne $\{1,3\}$ et $\{4,t\}$)
    \item $s$-$1$ (3) : ajouter $\to \{s,1,2,3,4,t\}$ (fusionne $\{s,2\}$ et $\{1,3,4,t\}$)
    \item $2$-$5$ (3) : ajouter $\to \{s,1,2,3,4,5,t\}$ (connecte $5$) $\checkmark$ \textbf{7 sommets, 6 ar\^etes = arbre complet}
\end{enumerate}

$$T = \{1\text{-}3(1),\ 4\text{-}t(1),\ s\text{-}2(2),\ 3\text{-}t(2),\ s\text{-}1(3),\ 2\text{-}5(3)\}$$
$$\text{Co\^ut MST} = 1+1+2+2+3+3 = \mathbf{12}$$

\medskip
\textbf{Q5. Puissance de chaque n\oe ud sur le MST :}

\begin{center}
\renewcommand{\arraystretch}{1.2}
\begin{tabular}{|c|l|c|}
\hline
\textbf{N\oe ud} & \textbf{Ar\^etes incidentes dans $T$} & $p(v)$ \\\hline
$s$ & $s$-$1$(3), $s$-$2$(2) & $\max(3,2) = 3$ \\\hline
$1$ & $s$-$1$(3), $1$-$3$(1) & $\max(3,1) = 3$ \\\hline
$2$ & $s$-$2$(2), $2$-$5$(3) & $\max(2,3) = 3$ \\\hline
$3$ & $1$-$3$(1), $3$-$t$(2) & $\max(1,2) = 2$ \\\hline
$4$ & $4$-$t$(1) & $1$ \\\hline
$5$ & $2$-$5$(3) & $3$ \\\hline
$t$ & $3$-$t$(2), $4$-$t$(1) & $\max(2,1) = 2$ \\\hline\hline
\multicolumn{2}{|r|}{\textbf{Puissance totale MST}} & $\mathbf{17}$ \\\hline
\end{tabular}
\end{center}

\medskip
\textbf{Q6. Algorithme BIP depuis $s$ :}

Initialisation : arbre $= \{s\}$, $p(s) = 0$.

\textbf{\'Etape 1 :} Ar\^etes depuis $s$ : $s$-$1$(3), $s$-$2$(2), $s$-$4$(5).

Co\^uts additionnels : $a(s\text{-}1) = 3-0=3$, $a(s\text{-}2) = 2-0=2$, $a(s\text{-}4) = 5-0=5$.

Min : $s$-$2$ (co\^ut $2$). $p(s)=2$, $p(2)=2$.

\textbf{\'Etape 2 :} Depuis $s$($p=2$), $2$($p=2$) :
\begin{itemize}[nosep]
    \item $s$-$1$(3) : $3-2=1$
    \item $s$-$4$(5) : $5-2=3$
    \item $2$-$4$(6) : $6-2=4$
    \item $2$-$5$(3) : $3-2=1$
\end{itemize}
Ex \ae quo : $s$-$1$(1) et $2$-$5$(1). Choisissons $s$-$1$. $p(s)=3$, $p(1)=3$.

\textbf{\'Etape 3 :} Depuis $s$($p=3$), $1$($p=3$), $2$($p=2$) :
\begin{itemize}[nosep]
    \item $1$-$3$(1) : $\max(1-3,0)=0$
    \item $1$-$4$(4) : $4-3=1$
    \item $2$-$5$(3) : $3-2=1$
    \item $s$-$4$(5) : $5-3=2$
\end{itemize}
Min : $1$-$3$ (co\^ut $0$). $p(1)$ inchang\'e, $p(3)=1$.

\textbf{\'Etape 4 :} Depuis $s$($p=3$), $1$($p=3$), $2$($p=2$), $3$($p=1$) :
\begin{itemize}[nosep]
    \item $3$-$t$(2) : $2-1=1$
    \item $3$-$4$(3) : $3-1=2$
    \item $1$-$4$(4) : $4-3=1$
    \item $2$-$5$(3) : $3-2=1$
\end{itemize}
Ex \ae quo : $3$-$t$(1), $1$-$4$(1), $2$-$5$(1). Choisissons $3$-$t$. $p(3)=2$, $p(t)=2$.

\textbf{\'Etape 5 :} Depuis $t$($p=2$) aussi :
\begin{itemize}[nosep]
    \item $4$-$t$(1) : $\max(1-2,0)=0$
    \item $1$-$4$(4) : $4-3=1$
    \item $2$-$5$(3) : $3-2=1$
    \item $5$-$t$(4) : $4-2=2$
\end{itemize}
Min : $4$-$t$ (co\^ut $0$). $p(t)$ inchang\'e, $p(4)=1$.

\textbf{\'Etape 6 :} Restant : $5$.
\begin{itemize}[nosep]
    \item $2$-$5$(3) : $3-2=1$
    \item $5$-$t$(4) : $4-2=2$
\end{itemize}
Min : $2$-$5$ (co\^ut $1$). $p(2)=3$, $p(5)=3$.

\medskip
Arbre BIP $= \{s\text{-}2(2),\ s\text{-}1(3),\ 1\text{-}3(1),\ 3\text{-}t(2),\ 4\text{-}t(1),\ 2\text{-}5(3)\}$

\begin{center}
\renewcommand{\arraystretch}{1.2}
\begin{tabular}{|c|c|c|c|c|c|c|c||c|}
\hline
\textbf{N\oe ud} & $s$ & $1$ & $2$ & $3$ & $4$ & $5$ & $t$ & \textbf{Total} \\\hline
$p_{\text{BIP}}$ & $3$ & $3$ & $3$ & $2$ & $1$ & $3$ & $2$ & $\mathbf{17}$ \\\hline
\end{tabular}
\end{center}

\medskip
\textbf{Q7. Comparaison :}

\begin{center}
\begin{tabular}{|l|c|c|}
\hline
& \textbf{MST (Kruskal)} & \textbf{BIP (depuis $s$)} \\\hline
Co\^ut de l'arbre & $12$ & $12$ \\\hline
Puissance totale & $17$ & $17$ \\\hline
\end{tabular}
\end{center}

Ici, les deux m\'ethodes donnent le \textbf{m\^eme arbre} et la m\^eme puissance totale. Ce n'est pas toujours le cas : BIP peut donner une puissance totale inf\'erieure au MST (cf. examens A et B).

\textbf{Remarque :} BIP est une heuristique qui ne garantit pas la puissance minimale, mais elle donne g\'en\'eralement un r\'esultat meilleur ou \'egal au MST en termes de puissance totale.
\end{correction}

\begin{attention}
\textbf{Strat\'egie d'examen --- D\'emarche syst\'ematique :}
\begin{enumerate}[nosep]
    \item \textbf{Identifier le type de graphe :} orient\'e/non orient\'e, poids n\'egatifs ?, circuits ?, DAG ?
    \item \textbf{Choisir l'algorithme :}
    \begin{itemize}[nosep]
        \item DAG $\to$ relaxation par niveaux (m\^eme si poids $< 0$)
        \item Poids $\geq 0$, pas DAG $\to$ Dijkstra
        \item Poids $< 0$, pas DAG $\to$ Bellman-Ford
    \end{itemize}
    \item \textbf{MST :} trier les ar\^etes, appliquer Kruskal, \textbf{v\'erifier les cycles}
    \item \textbf{BIP :} toujours calculer le co\^ut additionnel $= \max(C(u,v) - p(u), 0)$
    \item \textbf{Comparer :} puissance MST vs puissance BIP
\end{enumerate}
\end{attention}


\newpage
%=============================================================
%  SYNTHESE FINALE
%=============================================================
\section*{Synth\`ese : Quel algorithme pour quel probl\`eme ?}
\addcontentsline{toc}{section}{Synth\`ese}

\begin{center}
\renewcommand{\arraystretch}{1.5}
\begin{tabular}{|p{4.5cm}|p{4cm}|p{5cm}|}
\hline
\textbf{Probl\`eme} & \textbf{Algorithme} & \textbf{Condition / Complexit\'e} \\\hline
PCC, DAG (m\^eme poids $< 0$) & Relaxation topologique & $O(V+E)$. \textbf{Montrer les niveaux.} \\\hline
PCC, poids $\geq 0$ & Dijkstra & $O(E + V\log V)$. \textbf{Tableau d'ex\'ecution.} \\\hline
PCC, poids n\'egatifs, pas DAG & Bellman-Ford & $O(V \cdot E)$. D\'etecte circuits $-$. \\\hline
MST & Kruskal (ou Prim) & Trier + Union-Find. $O(E \log E)$. \\\hline
Puissance min sans fil & BIP & Co\^ut add. $= \max(C-p(u), 0)$. \\\hline
Flot maximum & Ford-Fulkerson / Edmonds-Karp & Graphe r\'esiduel + chemins augmentants. \\\hline
Chemins arcs-disjoints & Flot max, capacit\'es $= 1$ & Th. de Menger. \\\hline
Chemins sommets-disjoints & D\'edoublement + flot max & $v \to v_{in}, v_{out}$, cap. $= 1$. \\\hline
Capacit\'e nodale & D\'edoublement & Arc $v_{in} \to v_{out}$, cap. $= k$. \\\hline
Dominant connexe (CDS) & Glouton : degr\'e max & Voisin de D couvrant le plus. \\\hline
MPR & \'Etape oblig. + glouton & Voisinage 2 sauts. \\\hline
\end{tabular}
\end{center}

\vspace{2em}
\begin{center}
\textcolor{titrebleu}{\rule{0.6\textwidth}{1.5pt}}\\[0.5em]
\textit{\Large Bon courage pour l'examen !}\\[0.3em]
\textcolor{defvert}{\small Optimisation dans les R\'eseaux --- F.~Bendali-Mailfert --- S6 2025}
\end{center}

\end{document}
